\chapter{Sentiment Analysis}\label{sec-sentiment-analysis}

\textbf{Learning Objectives}

\emph{Overview starts}

\begin{itemize}
\tightlist
\item
  Understand what sentiment analysis is
\item
  How to analyze sentiment analysis
\item
  Learn how to understand sentiment over a timeline
\item
  Understand what to do with the emojis
\item
  How to use bigrams to better understand polarities
\item
  How to use word networks to extract polarized information
\item
  Correlate sentiment score and user ratings
\item
  Extract actionable information from topics, sentiment, and time
\item
  Classify emotions
\item
  Tackle prophanity in reviews
\end{itemize}

\emph{Overview ends}

\section{Meaning, Levels and Tasks of Sentiment
Analysis}\label{meaning-levels-and-tasks-of-sentiment-analysis}

Sentiment analysis is a widely used technique in text analytics. The
word ``sentiment'' is frequently interchanged with ``emotion,'' and it
has developed meanings in several subjects, such as literature, visual
arts, or psychology. Psychologists and philosophers often disagree about
their meaning. The word sentiment originates from the French term
``sentement'', and, as by the Oxford Dictionary, is ``What one feels
with regard to something; mental attitude (of approval or disapproval,
etc.); an opinion or view as to what is right or agreeable..'' In our
context, sentiment is an opinion expressed in a text, be it a review or
feedback that can be positive, negative, or neutral. Emotions may reveal
the expression of specific feelings such as happiness, joy, anger,
disgust, or surprise, and can be detected. Another name for sentiment
analysis is emotion detection in text corpora (Alslaity and Orji 2024).
In this chapter, we will analyze both sentiment and emotion. Sentiment
analysis is applied in several fields, such as finance Tirunillai and
Tellis (2012) to predict stock movements and assess economic indicators,
healthcare (Denecke and Deng 2015) to improve patient care, customer
service (K. S. Kumar, Desai, and Majumdar 2016) to improve customer
care, entertainment (Baid, Gupta, and Chaplot 2017) to understand
reactions to movies or music shows, product development (Fang and Zhan
2015) to improve product features, psychology (Xinyu Wang et al. 2013)
to analyze emotional states, politics (Bakliwal et al. 2013) to
understand opinion on policies or voter sentiments, social sciences
(Pozzi et al. 2016) to understand societal trends, or marketing Berger
et al. (2020) to analyze the valence of texts in social media posts,
consumer reviews, and other customer and firm-generated content.

In marketing, sentiment analysis can provide valuable insights into the
disaggregated and aggregated opinion about products and services over a
determined period of time on a specific topic or in comparison to
competing brands (see, e.g., T. Y. Lee and Bradlow (2011)). Based on the
language used, customers also can be classified by gender and age,
allowing for better market segmentation (see, e.g., Hartmann et al.
(2019)).

The corpus to be analyzed can be decomposed into conceptual levels. We
will use an example to understand better. Let us have a document, in our
case a short review: ``I purchased this tablet for Christmas and have
mixed feelings about it. The tablet quality was fantastic, but the WiFi
connection was disappointing.'' The highest level is the document level,
where sentiment analysis can be performed on each document of a
collection. In our review we have only one document. If the document is
a chapter, a page or a paragraph, its polarity can be assessed as
positive or negative. The next level is the sentence level. In some
instances, it is much better to consider a document at the sentence
level because the polarities of the sentences can be aggregated (see
Büschken and Allenby (2016)). Then comes the phrase level. In the
sentence, ``The tablet quality was fantastic, but the WiFi connection
was disappointing.'', the first phrase, ``The tablet quality was
fantastic,'' has a positive sentiment, while the second one, ``but the
WiFi connection was disappointing,'' has a negative sentiment.
Phrase-based sentiment analysis focuses on different parts of a
sentence, in contrast with sentence-based which yields an averaged
sentiment for the whole sentence. The last one is the aspect level.
Sentiment analysis focuses on specific aspects (topics) within a phrase.
For example, ``tablet quality'' (aspect) is associated with
``fantastic,'' which has a positive sentiment, while ``WiFi connection''
is associated with ``disappointing,'' which has a negative sentiment,
see Figure~\ref{fig-salevels}.

\begin{figure}

\centering{

\includegraphics[width=0.9\textwidth,height=\textheight]{figure/chapter8/salevels.png}

}

\caption{\label{fig-salevels}Sentiment Analysis Levels,
\newline Authors' own figure}

\end{figure}%

Information can be divided into facts and opinions. While facts are
expressed objectively, opinions are often representative of people's
biases, emotions, state of mind, or interests (Liu et al. 2010) and are
subjective. A document is deemed subjective if it contains adjectives,
not all of them but a subset of them (Rittman et al. 2004). Opinion
measurement is of great importance because it can influence companies'
or other people's behavior in very unpredictable ways (Atwal, Bryson,
and Kaiser 2022). Sentiment classification is one of the main tasks for
categorizing text into positive or negative, and the most common one is
polarity detection (Cambria et al. 2015). Cross-lingual classification
uses one language as training data to classify documents in another
language (Olsson, Oard, and Hajič 2005). Cross-domain classification is
similar, it uses a sentiment classifier trained on a source domain and
used on a different domain (Bollegala, Mu, and Goulermas 2015). Both
classifications are useful to cut annotation data costs. Deception
detection and opinion spam detection are processes focused on
identifying and classifying intentionally designed fake reviews aimed to
mislead and influence readers (Rosso and Cagnina 2017).\\
Sentiments can be expressed either explicitly, ``slow and unresponsive
touchscreen experience,'' or implicitly, ``I could use it for only 5
minutes''. Both sentences are an expression of a negative sentiment, and
there are several techniques such as subjectivity detection, polarity
detection, and sentiment strength detection (Saad 2014) that accomplish
this task. Aspects, or topics, can be identified using LDA (see Chapter
7); there are sentiment analysis aspect models that consider the
sentiment towards a specific topic included in the topic model (Lu et
al. 2011). There are several methods to accomplish this task, such as
syntax-based, supervised, unsupervised, or hybrid machine learning
(Schouten and Frasincar 2015). For a visual representation of the
sentiment analysis tasks, see Figure~\ref{fig-satasks}.

\begin{figure}

\centering{

\includegraphics[width=0.5\textwidth,height=\textheight]{figure/chapter8/satasks.png}

}

\caption{\label{fig-satasks}Sentiment Analysis Tasks, \newline Authors'
own figure}

\end{figure}%

\section{Sentiment Analysis
Methodology}\label{sentiment-analysis-methodology}

Several methods are available to perform sentiment analysis:
lexicon-based, machine learning-based, or hybrid approach Dhaoui,
Webster, and Tan (2017). An extra approach can be considered, too, that
contains transfer learning and aspect-based approaches (Wankhade, Rao,
and Kulkarni 2022). The lexicon-based approach is the most common one.
It uses a dictionary and deals with assigning scores to words. This can
be problematic, though, because, depending on the context, words can
have a positive or negative polarity, so domain detection is important
(H. Han et al. 2018). Such an algorithm takes a text, compares it
against a dictionary that has scores annotated with a specific polarity
and then aggregates them to form a final score (Taboada et al. 2011).
Given that we have studied classification (see Chapter 6), there are
several ways to accomplish text classification. What is needed is to
assign a score to a text, something that we frequently see on opinion
collections on the internet: a review connected to a score that ranges
normally between 1 and 5 stars. A hybrid method is a corpus-based
approach, which is a text classification problem that uses an annotated
data set based on a lexicon (Abdulla et al. 2013). Here, we disagree
with what is depicted in Figure~\ref{fig-saapproach} in (Wankhade, Rao,
and Kulkarni 2022), and we edited their image considering our
observations. Our figure presents various methods to perform sentiment
analysis, which are categorized into several approaches. Most of them
fall under the Machine Learning Approach, such as Linear Classifiers and
K-Nearest Neighbor; we do not refer to all of them here as they were
discussed in Chapter 6. Advanced techniques, such as Transformer-based
models, are discussed in Chapter 10, where the contextual relationships
in text data are explored.

\begin{figure}

\centering{

\includegraphics[width=0.9\textwidth,height=\textheight]{figure/chapter8/saapproach.png}

}

\caption{\label{fig-saapproach}Sentiment Analysis Approach}

\end{figure}%

\section{Performing Sentiment Analysis in
R}\label{performing-sentiment-analysis-in-r}

In R, several packages can be used to accomplish a sentiment analysis
task. \texttt{sentiment} (T. W. Rinker 2021), \texttt{syuzhet} (Matthew
L. Jockers 2015), and \texttt{SentimentAnalysis} (Proellochs and
Feuerriegel 2023) use the lexicon-based approach. They employ
dictionaries of words associated to a sentiment score.
\texttt{sentimentr}, uses the Jockers Rinker Sentiment Lexicon, based on
a combination of several lexicons. It takes into account several
interesting features, such as negators, amplifiers, and de-amplifiers,
to fine-tune the sentiment scores. The result is an intensity polarity
scale between very negative and very positive {[}-1,1{]}, passing
through a neutral 0. \texttt{syuzhet} uses scores from different
dictionaries, such as Bing (M. Hu and Liu 2004) that binarily classifies
words as positive or negative without an intensity, the NRC emotion
lexicon, Khoo and Johnkhan (2018), developed at the National Research
Council of Canada (NRC), which associates words to eight basic emotions,
such as anger, fear, anticipation, trust, surprise, joy and disgust, and
Afinn (Aung and Myo 2017), which assignes integer scores to words, that
discretely range from {[}-5, 5{]}. It also assigns positive or negative
sentiment to the result. \texttt{SentimentAnalysis} supports several
lexicons, such as the QDAP dictionary (T. Rinker, Rinker, and LazyData
2018), the Harvard IV dictionary (Frankel, Jennings, and Lee 2022), and
many more. \texttt{vader} (Roehrick 2020) implements the Valence Aware
Dictionary for Sentiment Reasoning (VADER, Hutto and Gilbert (2014))
which is a rule-based classifier. It uses a predefined dictionary and
considers the context of the words, being ``aware'' of the surrounding
words and their impact on the sentiment of the target word. Sentiment
scores are adjusted based on intensity modifiers, such as ``very'' or
``extremely'' or punctuation and capitalization (``GREAT!!'' is more
intense than ``great''), and negations are dealt with by flipping the
sentiment score.

The \texttt{sentiment.ai} package (Wiseman, Nydick, and Wisner 2022)
uses deep learning to accurately perform sentiment analysis on text that
is context-dependent. It uses transformer-based models, which are
pre-trained on large data sets, to identify the context in which sarcasm
and ambiguous phrases appear. Its multi-language support is a plus.
sentiment.ai seems to identify much better the negative context than
other packages \footnote{https://github.com/BenWiseman/sentiment.ai} as
by the authors' claims, being in disagreement with the lexicon-based
packages, while on the positive side there is agreement between them.

\subsection{Sentiment over a timeline}\label{sentiment-over-a-timeline}

Let us work with an example to better understand the steps to perform
sentiment analysis on a dataset. We use the same data we analyzed in the
previous chapters. Importing them into R and selecting the relevant
columns and rows should be a breeze for you by now. We chose two brands
because comparison is crucial in marketing analysis.

\begin{Shaded}
\begin{Highlighting}[]
\FunctionTok{library}\NormalTok{(dplyr)}
\FunctionTok{library}\NormalTok{(lubridate)}

\CommentTok{\# Prevent conversion of strings into factors}
\FunctionTok{options}\NormalTok{(}\AttributeTok{stringsAsFactors =} \ConstantTok{FALSE}\NormalTok{)}

\CommentTok{\# Read product reviews from the CSV file into a dataframe}
\NormalTok{prodreviews }\OtherTok{\textless{}{-}} \FunctionTok{read.csv}\NormalTok{(}\StringTok{"data/Product Reviews.csv"}\NormalTok{)}

\CommentTok{\# Select the desired columns}
\NormalTok{data }\OtherTok{\textless{}{-}}\NormalTok{ prodreviews[, }\FunctionTok{c}\NormalTok{(}\StringTok{"Content"}\NormalTok{, }\StringTok{"Rating"}\NormalTok{, }
                        \StringTok{"Brand"}\NormalTok{, }\StringTok{"Review\_date"}\NormalTok{)]}

\CommentTok{\# Transform the character content of Review\_date into date type}
\NormalTok{data}\SpecialCharTok{$}\NormalTok{Review\_date }\OtherTok{\textless{}{-}} \FunctionTok{mdy}\NormalTok{(data}\SpecialCharTok{$}\NormalTok{Review\_date)}

\CommentTok{\# Select the relevant rows (Amazon and Apple)}
\NormalTok{data }\OtherTok{\textless{}{-}} \FunctionTok{subset}\NormalTok{(data, Brand }\SpecialCharTok{\%in\%} \FunctionTok{c}\NormalTok{(}\StringTok{"Samsung"}\NormalTok{, }\StringTok{"Apple"}\NormalTok{, }\StringTok{"Amazon"}\NormalTok{))}
\end{Highlighting}
\end{Shaded}

After the dataset has been loaded, we must address an important detail:
emojis and emoticons. Text can contain them, sometimes in abundance, and
can influence the sentiment of the reviews and lead to better accuracy
of the overall scoreThese influence the sentiment of the reviews and
lead to better accuracy of the overall score, see Shelley Gupta, Singh,
and Ranjan (2020). There are also cases in which text can be
``emojified'', to better undertand feelings about certain products
(Hotz-Behofsits, Wlömert, and Nabout 2025). From the mentioned packages,
\texttt{vader} has a built-in lexicon that links emojis to sentiments.
There are proposed methods, though, that calculate sentiment analysis
from emojis by leveraging Tweets and GitHub posts containing emojis to
better classify text containing them (Z. Chen et al. 2021). In R, one
can use several packages to deal with emojis, such as \texttt{textclean}
where functions like \texttt{replace\_emoji()} or
\texttt{replace\_non\_ascii()}, transform emojis and emoticons into
their descriptive text counterpart. It is good to use them both, they
complete each other. Other such packages are \texttt{emoji}, or
\texttt{tidyEmoji}. We can also apply straightforward approaches as
follows: the identification of emojis and their transformation into text
descriptions. Because our dataset does not have emojis, the final result
will not be different from the initial text. The following
implementation is provided for your convenience so that you can use it
with your dataset.

\begin{Shaded}
\begin{Highlighting}[]
\FunctionTok{library}\NormalTok{(stringr)}

\CommentTok{\# Replace the emojis in reviews with their explanation}
\NormalTok{data}\SpecialCharTok{$}\NormalTok{Content }\OtherTok{\textless{}{-}} \FunctionTok{replace\_emoji}\NormalTok{(data}\SpecialCharTok{$}\NormalTok{Content)}

\CommentTok{\# Remove non{-}ASCII characters with suitable ASCII equivalent or removes them}
\NormalTok{data}\SpecialCharTok{$}\NormalTok{Content }\OtherTok{\textless{}{-}} \FunctionTok{replace\_non\_ascii}\NormalTok{(data}\SpecialCharTok{$}\NormalTok{Content)}
\end{Highlighting}
\end{Shaded}

After this procedure, the text is ready for analysis. In the
computation, we will include, several sentiment analysis packages, but
we will mostly focus on \texttt{sentiment.ai} in this chapter. It seems
that this package outperforms \texttt{sentimentAnalysis} or
\texttt{sentimentr} as by their author's claims on their GitHub page. We
keep the other results for comparison reasons.

\begin{Shaded}
\begin{Highlighting}[]
\FunctionTok{library}\NormalTok{(sentimentr)}
\FunctionTok{library}\NormalTok{(sentiment.ai)}
\FunctionTok{library}\NormalTok{(SentimentAnalysis)}
\FunctionTok{library}\NormalTok{(vader)}
\FunctionTok{library}\NormalTok{(syuzhet)}
\FunctionTok{library}\NormalTok{(data.table)}

\CommentTok{\# init\_sentiment.ai() is to be run once, at the beginning}

\CommentTok{\# sentiment.ai}
\NormalTok{sentiment.ai.score }\OtherTok{\textless{}{-}} \FunctionTok{sentiment\_score}\NormalTok{(data}\SpecialCharTok{$}\NormalTok{Content)}

\CommentTok{\# From Sentiment Analysis}
\NormalTok{sentimentAnalysis.score }\OtherTok{\textless{}{-}} \FunctionTok{analyzeSentiment}\NormalTok{(data}\SpecialCharTok{$}\NormalTok{Content)}\SpecialCharTok{$}\NormalTok{SentimentQDAP}

\CommentTok{\# From sentimentr}
\NormalTok{sentimentr.score }\OtherTok{\textless{}{-}} \FunctionTok{sentiment\_by}\NormalTok{(}\FunctionTok{get\_sentences}\NormalTok{(data}\SpecialCharTok{$}\NormalTok{Content),}
                    \DecValTok{1}\SpecialCharTok{:}\NormalTok{(}\FunctionTok{length}\NormalTok{(data}\SpecialCharTok{$}\NormalTok{Content)}\SpecialCharTok{{-}}\DecValTok{1}\NormalTok{))}\SpecialCharTok{$}\NormalTok{ave\_sentiment}
\CommentTok{\# Calculate and aggregate sentiment scores}

\CommentTok{\# From VADER}
\NormalTok{vader.score }\OtherTok{\textless{}{-}} \FunctionTok{vader\_df}\NormalTok{(data}\SpecialCharTok{$}\NormalTok{Content)}\SpecialCharTok{$}\NormalTok{compound}

\CommentTok{\# From syuzhet}
\NormalTok{syuzhet.score }\OtherTok{\textless{}{-}} \FunctionTok{get\_sentiment}\NormalTok{(data}\SpecialCharTok{$}\NormalTok{Content)}

\CommentTok{\# Compile results into a data table}
\NormalTok{comparison }\OtherTok{\textless{}{-}} \FunctionTok{data.table}\NormalTok{(}
  \AttributeTok{Review\_id =} \DecValTok{1}\SpecialCharTok{:}\FunctionTok{nrow}\NormalTok{(data),}
  \AttributeTok{target =}\NormalTok{ data}\SpecialCharTok{$}\NormalTok{Content,}
  \AttributeTok{sentiment.ai =}\NormalTok{ sentiment.ai.score,}
  \AttributeTok{sentimentAnalysis =}\NormalTok{ sentimentAnalysis.score,}
  \AttributeTok{sentimentr =}\NormalTok{ sentimentr.score,}
  \AttributeTok{vader =}\NormalTok{ vader.score,}
  \AttributeTok{syuzhet =}\NormalTok{ syuzhet.score, }
  \AttributeTok{Brand =}\NormalTok{ data}\SpecialCharTok{$}\NormalTok{Brand,}
  \AttributeTok{Rating =}\NormalTok{ data}\SpecialCharTok{$}\NormalTok{Rating,}
  \AttributeTok{R\_date =}\NormalTok{ data}\SpecialCharTok{$}\NormalTok{Review\_date}
\NormalTok{)}
\end{Highlighting}
\end{Shaded}

We transform the Date to display data on a monthly base and we average
\texttt{sentiment.ai} for each combiation of \texttt{Month} and
\texttt{Brand}. We have also renamed the column \texttt{Content} into
\texttt{target} for easier integration with other packages. The plot of
the graph follows next.

\begin{Shaded}
\begin{Highlighting}[]
\FunctionTok{library}\NormalTok{(ggplot2)}
\FunctionTok{library}\NormalTok{(lubridate)}

\CommentTok{\# Round \textquotesingle{}R\_date\textquotesingle{} down to the first day of its month, add it as a new column}
\NormalTok{comparison[, Month }\SpecialCharTok{:}\ErrorTok{=} \FunctionTok{floor\_date}\NormalTok{(R\_date, }\StringTok{"month"}\NormalTok{)]}

\CommentTok{\# Average score for selected columns}
\NormalTok{monthly\_aggregated }\OtherTok{\textless{}{-}}\NormalTok{ comparison[, .(}\AttributeTok{sentiment.ai =} \FunctionTok{mean}\NormalTok{(sentiment.ai)), }
\NormalTok{                          by }\OtherTok{=}\NormalTok{ .(Month, Brand)]}

\CommentTok{\# Plot the graph}
\FunctionTok{ggplot}\NormalTok{(monthly\_aggregated, }\FunctionTok{aes}\NormalTok{(}\AttributeTok{x =}\NormalTok{ Month, }\AttributeTok{y =}\NormalTok{ sentiment.ai, }
                               \AttributeTok{color =}\NormalTok{ Brand, }\AttributeTok{group =}\NormalTok{ Brand)) }\SpecialCharTok{+}
  \FunctionTok{geom\_smooth}\NormalTok{(}\AttributeTok{method =} \StringTok{"loess"}\NormalTok{, }\AttributeTok{se =} \ConstantTok{FALSE}\NormalTok{, }\AttributeTok{linewidth =} \FloatTok{1.2}\NormalTok{) }\SpecialCharTok{+} 
  \FunctionTok{geom\_point}\NormalTok{(}\AttributeTok{size =} \DecValTok{3}\NormalTok{) }\SpecialCharTok{+} \CommentTok{\# Slightly larger points}
  \FunctionTok{labs}\NormalTok{(}\AttributeTok{x =} \StringTok{"Month"}\NormalTok{,}
       \AttributeTok{y =} \StringTok{"Average sentiment.ai score"}\NormalTok{,}
       \AttributeTok{color =} \StringTok{"Brand"}
\NormalTok{      ) }\SpecialCharTok{+}
  \FunctionTok{theme\_minimal}\NormalTok{() }\SpecialCharTok{+}
  \FunctionTok{theme}\NormalTok{(}
    \AttributeTok{text =} \FunctionTok{element\_text}\NormalTok{(}\AttributeTok{size =} \DecValTok{12}\NormalTok{),}
    \AttributeTok{axis.text.x =} \FunctionTok{element\_text}\NormalTok{(}\AttributeTok{angle =} \DecValTok{45}\NormalTok{, }\AttributeTok{hjust =} \DecValTok{1}\NormalTok{)}
\NormalTok{  )}
\end{Highlighting}
\end{Shaded}

\begin{figure}[H]

\hfill\includegraphics{8.Chapter_files/figure-pdf/fig-sacomparison1a-1.pdf}

\caption{\label{fig-sacomparison1a}Monthly Aggregated sentiment.ai
Values Over Time for Each Brand \newline Authors' own image}

\end{figure}%

The graph in @fig-sacomparison1a shows, over a timeline, the monthly
aggregated sentiment scores elaborated with \texttt{sentiment.ai} for
the three selected brads: Amazon (represented mainly by Kindle Fire),
Apple and Samsung. The lines were smoothed using the LOESS method and
indicate a general sentiment over time. The Amazon line indicates a
positive trend, Apple's line apparently seems flat but it is constantly
increasing, with slight variations, while Samsung's trend line is more
volatile, indicating sentiment fluctuations. The Kindle Fire was
introduce later than its competitors, more precisely on November 15,
2011, nicely marketed before the Christmas holidays, showing it was
well-received by customers over time. Apple seems to have a stable and
increasing customer satisfaction with its iPad's while Samsung's
fluctuations in sentiment indicate mixed customer experiences, but the
trend seems overall positive.

A different type of visualization is provided in @fig-sacomparison2 to
show how sentiment towards the specified brands has evolved over time.
It highlights differences in how customer perception is captured by the
results of the ``sentiment.ai'' scores. This image provides more detail
for a month-to-month sentiment comparison. Amazon shows mostly positive
sentiment scores, hinting at generally favorable customer reception.
Apple exhibits a positive sentiment distribution, indicating a steady
customer satisfaction. Samsung has fluctuations in sentiment scores but
the overall sentiment, nevertheless, seems generally favorable.

\begin{Shaded}
\begin{Highlighting}[]
\FunctionTok{ggplot}\NormalTok{(monthly\_aggregated, }\FunctionTok{aes}\NormalTok{(}\AttributeTok{x =}\NormalTok{ Month, }\AttributeTok{y =}\NormalTok{ sentiment.ai, }
                               \AttributeTok{fill =}\NormalTok{ Brand)) }\SpecialCharTok{+}
  \FunctionTok{geom\_col}\NormalTok{(}\AttributeTok{show.legend =} \ConstantTok{FALSE}\NormalTok{) }\SpecialCharTok{+} \CommentTok{\# Use columns for visualization}
  \FunctionTok{labs}\NormalTok{(}
    \AttributeTok{x =} \StringTok{"Month"}\NormalTok{,}
    \AttributeTok{y =} \StringTok{"Average sentiment.ai score"}\NormalTok{,}
    \AttributeTok{fill =} \StringTok{"Brand"}
\NormalTok{  ) }\SpecialCharTok{+}
  \FunctionTok{theme\_minimal}\NormalTok{() }\SpecialCharTok{+}
  \FunctionTok{theme}\NormalTok{(}
    \AttributeTok{text =} \FunctionTok{element\_text}\NormalTok{(}\AttributeTok{size =} \DecValTok{12}\NormalTok{),}
    \AttributeTok{axis.text.x =} \FunctionTok{element\_text}\NormalTok{(}\AttributeTok{angle =} \DecValTok{45}\NormalTok{, }\AttributeTok{hjust =} \DecValTok{1}\NormalTok{)) }\SpecialCharTok{+}
  \FunctionTok{facet\_wrap}\NormalTok{(}\SpecialCharTok{\textasciitilde{}}\NormalTok{ Brand, }\AttributeTok{ncol =} \DecValTok{1}\NormalTok{, }\AttributeTok{scales =} \StringTok{"free\_y"}\NormalTok{)}
\end{Highlighting}
\end{Shaded}

\begin{figure}[H]

\centering{

\includegraphics{8.Chapter_files/figure-pdf/fig-sacomparison2-1.pdf}

}

\caption{\label{fig-sacomparison2}Monthly Aggregated Sentiment.ai Values
for Each Brand \newline Authors' own figure}

\end{figure}%

\subsection{Sentiment word networks}\label{sentiment-word-networks}

Now, let us gain a better understanding of Samsung's positive and
negative feedback through the use of bigrams and word networks. Before
we dive into details, bigrams are pairs of consecutive words that
represent common word associations used in context. Word networks
represent the relationships between words represented as graphs. The
nodes correspond to words, the edges indicate their co-occurrence. Both
make it easier to identify key themes or structures in textual data. To
illustrate this process we select the Amazon value from the Brand column
and separate them into strongly positive and negative groups, where the
score is either more positive than 0.7, or more negative than -0.7. Then
we do the preprocessing part, and we set up the dataset for
visualizations.

\begin{Shaded}
\begin{Highlighting}[]
\FunctionTok{library}\NormalTok{(dplyr)}
\FunctionTok{library}\NormalTok{(tm)}
\FunctionTok{library}\NormalTok{(SnowballC)}
\FunctionTok{library}\NormalTok{(wordcloud)}
\FunctionTok{library}\NormalTok{(slam)}
\FunctionTok{library}\NormalTok{(RWeka)}
\FunctionTok{library}\NormalTok{(tidyr)}
\FunctionTok{library}\NormalTok{(igraph)}
\FunctionTok{library}\NormalTok{(ggraph)}

\CommentTok{\# Filter Amazon reviews and split into positive and negative}
\NormalTok{kindle\_reviews }\OtherTok{\textless{}{-}}\NormalTok{ comparison }\SpecialCharTok{\%\textgreater{}\%} \FunctionTok{filter}\NormalTok{(Brand }\SpecialCharTok{==} \StringTok{"Amazon"}\NormalTok{)}
\NormalTok{pos\_reviews }\OtherTok{\textless{}{-}}\NormalTok{ kindle\_reviews }\SpecialCharTok{\%\textgreater{}\%} \FunctionTok{filter}\NormalTok{(sentiment.ai }\SpecialCharTok{\textgreater{}} \FloatTok{0.7}\NormalTok{)}
\NormalTok{neg\_reviews }\OtherTok{\textless{}{-}}\NormalTok{ kindle\_reviews }\SpecialCharTok{\%\textgreater{}\%} \FunctionTok{filter}\NormalTok{(sentiment.ai }\SpecialCharTok{\textless{}} \SpecialCharTok{{-}}\FloatTok{0.7}\NormalTok{)}

\CommentTok{\# Text preprocessing function}
\NormalTok{preprocess }\OtherTok{\textless{}{-}} \ControlFlowTok{function}\NormalTok{(text) \{}
\NormalTok{  text }\SpecialCharTok{\%\textgreater{}\%} \FunctionTok{tolower}\NormalTok{() }\SpecialCharTok{\%\textgreater{}\%} 
    \FunctionTok{removeWords}\NormalTok{(}\FunctionTok{stopwords}\NormalTok{(}\StringTok{"english"}\NormalTok{)) }\SpecialCharTok{\%\textgreater{}\%} 
    \FunctionTok{removePunctuation}\NormalTok{(}\AttributeTok{ucp =} \ConstantTok{TRUE}\NormalTok{) }\SpecialCharTok{\%\textgreater{}\%}
    \FunctionTok{stripWhitespace}\NormalTok{() }\SpecialCharTok{\%\textgreater{}\%}
    \FunctionTok{wordStem}\NormalTok{(}\AttributeTok{language =} \StringTok{"en"}\NormalTok{) }\SpecialCharTok{\%\textgreater{}\%}
    \FunctionTok{removeNumbers}\NormalTok{()}
\NormalTok{\}}

\CommentTok{\# Preprocess text and transform into a corpus}
\NormalTok{corpus\_pos }\OtherTok{\textless{}{-}} \FunctionTok{VCorpus}\NormalTok{(}\FunctionTok{VectorSource}\NormalTok{(}\FunctionTok{lapply}\NormalTok{(pos\_reviews}\SpecialCharTok{$}\NormalTok{target, preprocess)))}
\NormalTok{corpus\_neg }\OtherTok{\textless{}{-}} \FunctionTok{VCorpus}\NormalTok{(}\FunctionTok{VectorSource}\NormalTok{(}\FunctionTok{lapply}\NormalTok{(neg\_reviews}\SpecialCharTok{$}\NormalTok{target, preprocess)))}

\CommentTok{\# Bigram function}
\NormalTok{bigram\_tokenizer }\OtherTok{\textless{}{-}} \ControlFlowTok{function}\NormalTok{(x) }\FunctionTok{NGramTokenizer}\NormalTok{(x, }
                                               \FunctionTok{Weka\_control}\NormalTok{(}\AttributeTok{min =} \DecValTok{2}\NormalTok{, }\AttributeTok{max =} \DecValTok{2}\NormalTok{))}

\CommentTok{\# Create matrices for both positive and negative reviews}
\NormalTok{process\_bigrams }\OtherTok{\textless{}{-}} \ControlFlowTok{function}\NormalTok{(corpus) \{}
\NormalTok{  tdm }\OtherTok{\textless{}{-}} \FunctionTok{TermDocumentMatrix}\NormalTok{(corpus, }
                            \AttributeTok{control =} \FunctionTok{list}\NormalTok{(}\AttributeTok{tokenize =}\NormalTok{ bigram\_tokenizer))}
  \FunctionTok{removeSparseTerms}\NormalTok{(tdm, }\FloatTok{0.999}\NormalTok{) }\SpecialCharTok{\%\textgreater{}\%} 
  \FunctionTok{as.matrix}\NormalTok{() }\SpecialCharTok{\%\textgreater{}\%} 
  \FunctionTok{rowSums}\NormalTok{() }\SpecialCharTok{\%\textgreater{}\%}
  \FunctionTok{sort}\NormalTok{(}\AttributeTok{decreasing =} \ConstantTok{TRUE}\NormalTok{)}
\NormalTok{\}}

\NormalTok{v\_pos\_bigram }\OtherTok{\textless{}{-}} \FunctionTok{process\_bigrams}\NormalTok{(corpus\_pos)}
\NormalTok{v\_neg\_bigram }\OtherTok{\textless{}{-}} \FunctionTok{process\_bigrams}\NormalTok{(corpus\_neg)}

\CommentTok{\# Create data frames}
\NormalTok{df\_pos\_bigram }\OtherTok{\textless{}{-}} \FunctionTok{data.frame}\NormalTok{(}\AttributeTok{word =} \FunctionTok{names}\NormalTok{(v\_pos\_bigram), }
                            \AttributeTok{freq =}\NormalTok{ v\_pos\_bigram)}
\NormalTok{df\_neg\_bigram }\OtherTok{\textless{}{-}} \FunctionTok{data.frame}\NormalTok{(}\AttributeTok{word =} \FunctionTok{names}\NormalTok{(v\_neg\_bigram), }
                            \AttributeTok{freq =}\NormalTok{ v\_neg\_bigram)}
\end{Highlighting}
\end{Shaded}

So far, only preprocessing has been done. What we are doing now is
creating a function for presenting bigram word frequencies, the positive
and negative ones, one on top of the other, so that we can compare them.

\begin{Shaded}
\begin{Highlighting}[]
\CommentTok{\# Plot barplots}
\NormalTok{plot\_word\_data }\OtherTok{\textless{}{-}} \ControlFlowTok{function}\NormalTok{(d, title) \{}
  \FunctionTok{par}\NormalTok{(}\AttributeTok{mar =} \FunctionTok{c}\NormalTok{(}\DecValTok{5}\NormalTok{, }\DecValTok{4}\NormalTok{, }\DecValTok{4}\NormalTok{, }\DecValTok{2}\NormalTok{) }\SpecialCharTok{+} \FloatTok{0.1}\NormalTok{)  }\CommentTok{\# Adjust margins}
  \FunctionTok{barplot}\NormalTok{(d[}\DecValTok{1}\SpecialCharTok{:}\DecValTok{30}\NormalTok{, ]}\SpecialCharTok{$}\NormalTok{freq, }\AttributeTok{las =} \DecValTok{2}\NormalTok{, }\AttributeTok{names.arg =}\NormalTok{ d[}\DecValTok{1}\SpecialCharTok{:}\DecValTok{30}\NormalTok{, ]}\SpecialCharTok{$}\NormalTok{word, }
          \AttributeTok{col =} \StringTok{"lightblue"}\NormalTok{, }\AttributeTok{main =}\NormalTok{ title, }\AttributeTok{ylab =} \StringTok{"Word frequencies"}\NormalTok{, }
          \AttributeTok{cex.names =} \FloatTok{0.6}\NormalTok{)}
\NormalTok{\}}

\CommentTok{\# Plot the positive and negative reviews one under the other}
\FunctionTok{par}\NormalTok{(}\AttributeTok{mfrow =} \FunctionTok{c}\NormalTok{(}\DecValTok{2}\NormalTok{, }\DecValTok{1}\NormalTok{))}
\FunctionTok{plot\_word\_data}\NormalTok{(df\_pos\_bigram, }\StringTok{"Most frequent bigrams in positive reviews"}\NormalTok{)}
\FunctionTok{plot\_word\_data}\NormalTok{(df\_neg\_bigram, }\StringTok{"Most frequent bigrams in negative reviews"}\NormalTok{)}
\end{Highlighting}
\end{Shaded}

\begin{figure}[H]

\centering{

\includegraphics{8.Chapter_files/figure-pdf/fig-posneg2-1.pdf}

}

\caption{\label{fig-posneg2}}

\end{figure}%

We can notice from the plots in Figure~\ref{fig-posneg2} containing the
first 30 most frequent bigrams, positive and negative reviews contain
``kindle fire'' but, after that, they start diverging. Positive reviews
focus on customers' satisfaction with the product's features, such as
ease of use, emotional attachment to it, and the suitability as a gift.
Negative reviews unveil customer service issues, several technical
problems, and dissatisfaction with some features or performance of the
device. The comparison with iPad seems to be considered in negative
reviews, perhaps, by considering Apple devices a better deal.

\begin{Shaded}
\begin{Highlighting}[]
\CommentTok{\# Hypothetical frequency sums based on visual inspection:}
\NormalTok{positive\_total }\OtherTok{\textless{}{-}} \FunctionTok{sum}\NormalTok{(df\_pos\_bigram}\SpecialCharTok{$}\NormalTok{freq[}\DecValTok{1}\SpecialCharTok{:}\DecValTok{30}\NormalTok{])}
\NormalTok{negative\_total }\OtherTok{\textless{}{-}} \FunctionTok{sum}\NormalTok{(df\_neg\_bigram}\SpecialCharTok{$}\NormalTok{freq[}\DecValTok{1}\SpecialCharTok{:}\DecValTok{30}\NormalTok{])}

\FunctionTok{print}\NormalTok{(positive\_total)}
\end{Highlighting}
\end{Shaded}

\begin{verbatim}
[1] 3677
\end{verbatim}

\begin{Shaded}
\begin{Highlighting}[]
\FunctionTok{print}\NormalTok{(negative\_total)}
\end{Highlighting}
\end{Shaded}

\begin{verbatim}
[1] 1386
\end{verbatim}

In the above frequency count, we notice that the most frequent bigrams
show about three times higher positive counts than the negative ones,
but the situation changes for the total number of bigrams which we
retrieve next.

\begin{Shaded}
\begin{Highlighting}[]
\CommentTok{\# Sum the frequencies of positive/negative bigrams}
\NormalTok{positive\_total }\OtherTok{\textless{}{-}} \FunctionTok{sum}\NormalTok{(df\_pos\_bigram}\SpecialCharTok{$}\NormalTok{freq)}
\NormalTok{negative\_total }\OtherTok{\textless{}{-}} \FunctionTok{sum}\NormalTok{(df\_neg\_bigram}\SpecialCharTok{$}\NormalTok{freq)}

\FunctionTok{print}\NormalTok{(}\FunctionTok{paste}\NormalTok{(}\StringTok{"Total frequency of positive bigrams:"}\NormalTok{, positive\_total))}
\end{Highlighting}
\end{Shaded}

\begin{verbatim}
[1] "Total frequency of positive bigrams: 18693"
\end{verbatim}

\begin{Shaded}
\begin{Highlighting}[]
\FunctionTok{print}\NormalTok{(}\FunctionTok{paste}\NormalTok{(}\StringTok{"Total frequency of negative bigrams:"}\NormalTok{, negative\_total))}
\end{Highlighting}
\end{Shaded}

\begin{verbatim}
[1] "Total frequency of negative bigrams: 12901"
\end{verbatim}

A word network would give further insight into their connection and
frequency. Let us plot the positive terms and better understand the
salient features that make customers happy about a brand. The creation
of a function to be used in a versatile manner for both positive and
negative reviews is a handy tool and it is presented in the code below.

\begin{Shaded}
\begin{Highlighting}[]
\CommentTok{\# Build and visualize word networks}
\NormalTok{visualize\_word\_network }\OtherTok{\textless{}{-}} \ControlFlowTok{function}\NormalTok{(d, freq\_threshold, title) \{}
\NormalTok{  d\_sep }\OtherTok{\textless{}{-}}\NormalTok{ d }\SpecialCharTok{\%\textgreater{}\%} 
    \FunctionTok{separate}\NormalTok{(word, }\FunctionTok{c}\NormalTok{(}\StringTok{"word1"}\NormalTok{, }\StringTok{"word2"}\NormalTok{), }\AttributeTok{sep =} \StringTok{" "}\NormalTok{)}

\NormalTok{wordnet }\OtherTok{\textless{}{-}}\NormalTok{ d\_sep }\SpecialCharTok{\%\textgreater{}\%} 
    \FunctionTok{filter}\NormalTok{(freq }\SpecialCharTok{\textgreater{}}\NormalTok{ freq\_threshold) }\SpecialCharTok{\%\textgreater{}\%} 
    \FunctionTok{graph\_from\_data\_frame}\NormalTok{()}
  
\FunctionTok{ggraph}\NormalTok{(wordnet, }\AttributeTok{layout =} \StringTok{"fr"}\NormalTok{) }\SpecialCharTok{+}
    \FunctionTok{geom\_edge\_link}\NormalTok{(}\FunctionTok{aes}\NormalTok{(}\AttributeTok{color =}\NormalTok{ freq, }\AttributeTok{width =}\NormalTok{ freq), }
        \AttributeTok{arrow =}\NormalTok{ grid}\SpecialCharTok{::}\FunctionTok{arrow}\NormalTok{(}\AttributeTok{type =} \StringTok{"closed"}\NormalTok{, }\AttributeTok{length =} \FunctionTok{unit}\NormalTok{(}\FloatTok{0.15}\NormalTok{, }\StringTok{"cm"}\NormalTok{))) }\SpecialCharTok{+}
    \FunctionTok{geom\_node\_point}\NormalTok{() }\SpecialCharTok{+}
    \FunctionTok{geom\_node\_text}\NormalTok{(}\FunctionTok{aes}\NormalTok{(}\AttributeTok{label =}\NormalTok{ name), }\AttributeTok{repel =} \ConstantTok{TRUE}\NormalTok{, }\AttributeTok{size =} \DecValTok{3}\NormalTok{) }\SpecialCharTok{+}
    \FunctionTok{scale\_edge\_color\_gradient}\NormalTok{(}\AttributeTok{low =} \StringTok{"lightblue"}\NormalTok{, }\AttributeTok{high =} \StringTok{"blue"}\NormalTok{) }\SpecialCharTok{+}
    \FunctionTok{scale\_edge\_width}\NormalTok{(}\AttributeTok{range =} \FunctionTok{c}\NormalTok{(}\FloatTok{0.5}\NormalTok{, }\DecValTok{3}\NormalTok{)) }\SpecialCharTok{+}
    \FunctionTok{theme\_void}\NormalTok{() }\SpecialCharTok{+}
    \FunctionTok{labs}\NormalTok{(}\AttributeTok{title =}\NormalTok{ title, }\AttributeTok{edge\_color =} \StringTok{"Frequency"}\NormalTok{, }\AttributeTok{edge\_width =} \StringTok{"Frequency"}\NormalTok{)}
\NormalTok{\}}

\CommentTok{\# Visualize the word network for positive reviews}
\FunctionTok{visualize\_word\_network}\NormalTok{(df\_pos\_bigram, }\DecValTok{30}\NormalTok{, }\StringTok{"Word Network {-} Positive Reviews"}\NormalTok{)}
\end{Highlighting}
\end{Shaded}

\includegraphics{8.Chapter_files/figure-pdf/posneg5-1.pdf}

In the above coding example and the accompanying graph that results from
it, Kindle Fire, or simply Fire, is the main product. The central node
is ``kindle'' and is strongly connected to ``fire'', indicating this as
the most frequently bigram in these reviews. This is as expected due to
the product's name and so far no new information. But ``kindle'' is
associated with other positive terms, which frequency are emphasized by
a thicker connection line between them. Connections such as ``purchased
kindle'', ``new kindle'', ``received kindle'', or ``bought kindle'',
show that many reviewers talk about the buying experience. The user
experience expressed by terms such as ``easy use'', or ``set easy'',
indicate the users find the item easy to use and set up. Positive
sentiments are expressed through ``absolutely love'', ``really like'',
and are strong selling points for the Amazon brand. Words like ``great
product'' and ``works great'' indicate pleasance with the product's
performance. Gift and holiday mentions, such as ``christmas gift,''best
gift,``, and''present'', indicate that the tablets indicate a high
propensity to sell these products around main holiday events. Multimedia
capabilities are expressed through bigrams such as ``watch movies,'',
``play/playing games.'' A good selling point is expressed by ``read(ing)
books,'' and ``can read''. A positive brand association is done through
words like ``thank(s) amazon'', showing appreciation for the brand and
the product. User friendliness, checking emails, product pricing and
overall satisfaction are reflected through ``friendly user,'', ``email
check,'', ``price product,'', or ``one best.'' But let us analyze the
negative reviews as well.

\begin{Shaded}
\begin{Highlighting}[]
\CommentTok{\# Visualize the word network for positive reviews}
\FunctionTok{visualize\_word\_network}\NormalTok{(df\_neg\_bigram, }\DecValTok{10}\NormalTok{, }\StringTok{"Word Network {-} Negative Reviews"}\NormalTok{)}
\end{Highlighting}
\end{Shaded}

\includegraphics{8.Chapter_files/figure-pdf/posneg6-1.pdf}

Obviously, the central term is again ``kindle'', combined with ``fire'',
which is the central focus of many negative reviews too. This is
strongly connected to ``disappointed'', indicated a bad user experience.
The purchase experience is discussed too, through terms like ``purchased
kindle,'' ``new kindle,'' ``received kindle,'' ``bought kindle,'' this
time in a negative context. A major dissatisfaction is given by the
``customer service,'' or ``customer support''. Technical problems and
malfunctionings are indicated through ``power button,'' ``stopped
working,'', ``tech support,'', ``wifi connection,'', ``connect wifi,'',
``screen touch,'', ``control volume,'', ``battery life,'' ``slow
loading,'' or ``freezing''. Strong negative sentiments are expressed
through harsh terms, such as ``junk piece,''piece crap,'' or with
elements of frustration, ``waste money,'', ``every time,'', ``long
time.'' A painful touch for a company is when their products are
unfavorably compared to the competitors' products, considering them
better, ``buy ipad'', ``got ipad'', expressing dissastifaction. Other
important terms that reflect other issues, are ``credit card,'' ``called
support,'' ``membership store,'' or ``outside us'', reflecting payment
problems or regional restrictions. Interestingly, the network of
negative reviews is much denser that the previously discussed network of
positive reviews. Detecting and managing negative feedback is extremely
important, if ignored, it can severely damage the brand's reputation as
companies need to keep their customer base and/or trying to acquire new
customers.

\subsection{Negative text sections
detection}\label{negative-text-sections-detection}

For a broader view and detecting exactly which sentences are negative in
a review, there are several possibilities to implement that. A quick
solution can be achieved with the help of the \texttt{sentiment.ai}
package in conjunction with the \texttt{sentimentr} package.

\begin{Shaded}
\begin{Highlighting}[]
\FunctionTok{library}\NormalTok{(magrittr)}
\FunctionTok{library}\NormalTok{(purrr)}
\CommentTok{\# Set seed for reproducibility}
\FunctionTok{set.seed}\NormalTok{(}\DecValTok{42}\NormalTok{)}

\CommentTok{\# Subdivision into sentences and score assignment to each of them  }
\CommentTok{\# Filter data and sample one row per brand}
\NormalTok{output }\OtherTok{\textless{}{-}}\NormalTok{ comparison }\SpecialCharTok{\%\textgreater{}\%}
  \FunctionTok{filter}\NormalTok{(Brand }\SpecialCharTok{\%in\%} \FunctionTok{c}\NormalTok{(}\StringTok{"Apple"}\NormalTok{, }\StringTok{"Samsung"}\NormalTok{)) }\SpecialCharTok{\%\textgreater{}\%}
  \FunctionTok{filter}\NormalTok{(sentiment.ai }\SpecialCharTok{\textless{}} \SpecialCharTok{{-}}\FloatTok{0.6}\NormalTok{) }\SpecialCharTok{\%\textgreater{}\%}
  \FunctionTok{group\_by}\NormalTok{(Brand) }\SpecialCharTok{\%\textgreater{}\%}
  \FunctionTok{slice\_sample}\NormalTok{(}\AttributeTok{n =} \DecValTok{1}\NormalTok{) }\SpecialCharTok{\%\textgreater{}\%}
  \FunctionTok{ungroup}\NormalTok{() }\SpecialCharTok{\%\textgreater{}\%}
  \CommentTok{\# Subdivide reviews into sentences}
  \FunctionTok{mutate}\NormalTok{(}\AttributeTok{review =} \FunctionTok{map}\NormalTok{(target, get\_sentences)) }\SpecialCharTok{\%\textgreater{}\%}
  \FunctionTok{unnest}\NormalTok{(review)}

\CommentTok{\# Calculate sentiment scores by review}
\NormalTok{sentiment\_output }\OtherTok{\textless{}{-}} \FunctionTok{sentiment\_by}\NormalTok{(output}\SpecialCharTok{$}\NormalTok{review, output}\SpecialCharTok{$}\NormalTok{Review\_id)}

\CommentTok{\# Highlight sentiment results}
\FunctionTok{highlight}\NormalTok{(sentiment\_output)}
\end{Highlighting}
\end{Shaded}

\begin{figure}[H]

{\centering \includegraphics[width=1\textwidth,height=\textheight]{figure/chapter8/smr.png}

}

\caption{Sentence-based Sentiment Analysis}

\end{figure}%

We select the rows by the \texttt{sentiment.ai} score because the
algorithm performs better at the review level. We chose a higher
negative intensity (-0.6), which means that we are first narrowing down
heavily to reviews with an overall negative sentiment. A threshold of
-0.6 is a practical and empirical (rather than absolute) value chosen by
us to keep a balance between ensuring that the most negative reviews are
selected but at the same time we keep enough data for analysis.
\texttt{sentimentr} performs at the sentence level. By using the
\texttt{sentiment.ai} package at document level first, we select the
strongly negative reviews, then, at the second step, we use
\texttt{sentimentr} to identify specific sentences that drive that
negativity. This example shows that by pinpointing the overall negative
results, we might obtain more precise results when we act at the
sentence level. Both reviews are randomly selected, but we control the
randomness with the set.seed() function for reproducibility reasons. The
first review is of an Apple product, and the second is a Samsung.

\subsection{Sentiment Score and
Ratings}\label{sentiment-score-and-ratings}

It is extremely frequent to be asked to encode an opinion on a product
or service into a numerical value, typically a ``star'' value between 1
and 5, with one being the most negative and five the most positive. This
representation is a quick snapshot of the overall satisfaction given by
a customer, and it is a categorical variable (Al-Natour and Turetken
2020). One can understand that this variable is not ordinal because the
subjective opinion of each reviewer gives it, and, as such, the
``distance'' between 1 star and 2 is not the same as the distance
between 4 and 5. Is this reflected in the sentiment score, too?
Sometimes yes, sometimes not. The emotions expressed in words can travel
on a different path than what is given as a rating. Combining these two
aspects, we might get a more comprehensive understanding of customer
feedback. Is a sentiment score directly correlated to the rating given?
Let us discover it with a scatter plot.

In many marketing contexts, customers, when they write a review, also
encode their opinion about a product or service into a numerical value,
typically a ``star'' value between 1 and 5, with one being the most
negative and five the most positive. This representation is a quick
snapshot of the overall satisfaction given by a customer, and it is a
categorical variable (Al-Natour and Turetken 2020). One can understand
that this variable is not ordinal because the subjective interpretation
of each reviewer of the ``distance'' between 1 star and 2 is not
necessarily the same as the distance between 4 and 5. In fact, in
real-world settings, the majority of customers typically contribute
overly favorable ratings, leading on the aggregate to so-called J-shape
distributions. Such rating distributions are extremely skewed with a
mass at the positive end of the rating scale (e.g., Moe, Netzer, and
Schweidel (2017)). Similarly, the emotions expressed in words can travel
on a different path than what is given as a rating. The automatically
calculated sentiment score, tries to reveal an emotional content that
star ratings alone may not capture. Combining these two aspects, we
might get a more comprehensive understanding of customer feedback. Is a
sentiment score directly correlated to the rating given? The idea is to
compare these two metrics to see if meaningful differences emerge. Let
us discover it with a scatter plot.

\begin{Shaded}
\begin{Highlighting}[]
\CommentTok{\# Boxplot to visualize the relationship between sentiment.ai scores and Rating}
\FunctionTok{ggplot}\NormalTok{(comparison, }\FunctionTok{aes}\NormalTok{(}\AttributeTok{x =} \FunctionTok{factor}\NormalTok{(Rating), }\AttributeTok{y =}\NormalTok{ sentiment.ai)) }\SpecialCharTok{+}
  \FunctionTok{geom\_boxplot}\NormalTok{() }\SpecialCharTok{+}
  \FunctionTok{labs}\NormalTok{(}
    \AttributeTok{title =} \StringTok{"Distribution of Sentiment Scores by Star Rating"}\NormalTok{,}
    \AttributeTok{x =} \StringTok{"Star Rating"}\NormalTok{,}
    \AttributeTok{y =} \StringTok{"Sentiment.ai Score"}
\NormalTok{  ) }\SpecialCharTok{+}
  \FunctionTok{theme\_minimal}\NormalTok{()}
\end{Highlighting}
\end{Shaded}

\begin{figure}[H]

\centering{

\includegraphics{8.Chapter_files/figure-pdf/fig-sentrating1-1.pdf}

}

\caption{\label{fig-sentrating1}}

\end{figure}%

In Figure~\ref{fig-sentrating1}, each box represents the distribution of
sentiment scores for a given star rating. The horizontal line within
each box is the median sentiment score, and the upper and lower edges
indicate the quartiles. Any dots outside the whiskers show outliers. It
can be noticed that there are reviews with higher sentiment scores that
occur at lower star ratings, highlighting that emotional content may not
strictly align with the numeric rating. Higher sentiment scores
generally correspond to higher ratings and the opposite for lower
values. Is the correlation between these two values strong?

\begin{Shaded}
\begin{Highlighting}[]
\CommentTok{\# Calculate the correlation between sentiment.ai and Rating}
\NormalTok{correlation }\OtherTok{\textless{}{-}} \FunctionTok{cor}\NormalTok{(comparison}\SpecialCharTok{$}\NormalTok{sentiment.ai, comparison}\SpecialCharTok{$}\NormalTok{Rating, }\AttributeTok{use =} \StringTok{"complete.obs"}\NormalTok{)}
\FunctionTok{cat}\NormalTok{(}\StringTok{"Correlation between sentiment.ai and Rating:"}\NormalTok{, correlation, }\StringTok{"}\SpecialCharTok{\textbackslash{}n}\StringTok{"}\NormalTok{)}
\end{Highlighting}
\end{Shaded}

\begin{verbatim}
Correlation between sentiment.ai and Rating: 0.6251258 
\end{verbatim}

The results indicate a moderate to strong correlation, meaning that the
sentiment score can be a good predictor of ratings. Let us check
visually if these correlations match.

\begin{Shaded}
\begin{Highlighting}[]
\FunctionTok{set.seed}\NormalTok{(}\DecValTok{42}\NormalTok{) }\CommentTok{\# For reproducibility}

\CommentTok{\# Examples of reviews with high sentiment.ai and high rating}
\NormalTok{high\_sentiment\_high\_rating }\OtherTok{\textless{}{-}}\NormalTok{ comparison }\SpecialCharTok{\%\textgreater{}\%}
  \FunctionTok{filter}\NormalTok{(sentiment.ai }\SpecialCharTok{\textgreater{}} \FloatTok{0.5}\NormalTok{, Rating }\SpecialCharTok{==} \DecValTok{5}\NormalTok{) }\SpecialCharTok{\%\textgreater{}\%}
  \FunctionTok{slice\_sample}\NormalTok{(}\AttributeTok{n =} \DecValTok{3}\NormalTok{)}
\FunctionTok{cat}\NormalTok{(}\StringTok{"High Sentiment, High Rating Reviews:}\SpecialCharTok{\textbackslash{}n}\StringTok{"}\NormalTok{)}
\end{Highlighting}
\end{Shaded}

\begin{verbatim}
High Sentiment, High Rating Reviews:
\end{verbatim}

\begin{Shaded}
\begin{Highlighting}[]
\CommentTok{\# Display values in a nicely formatted table}
\FunctionTok{kable}\NormalTok{(high\_sentiment\_high\_rating[, .(Review\_id, target, sentiment.ai, Rating)], }
      \AttributeTok{format=}\StringTok{"markdown"}\NormalTok{) }\SpecialCharTok{\%\textgreater{}\%}
  \FunctionTok{kable\_styling}\NormalTok{(}\AttributeTok{full\_width =} \ConstantTok{FALSE}\NormalTok{) }\SpecialCharTok{\%\textgreater{}\%}
  \FunctionTok{column\_spec}\NormalTok{(}\DecValTok{1}\NormalTok{, }\AttributeTok{width =} \StringTok{"4em"}\NormalTok{) }\SpecialCharTok{\%\textgreater{}\%}   
  \FunctionTok{column\_spec}\NormalTok{(}\DecValTok{2}\NormalTok{, }\AttributeTok{width =} \StringTok{"25em"}\NormalTok{) }\SpecialCharTok{\%\textgreater{}\%}    
  \FunctionTok{column\_spec}\NormalTok{(}\DecValTok{3}\NormalTok{, }\AttributeTok{width =} \StringTok{"5em"}\NormalTok{)  }
\end{Highlighting}
\end{Shaded}

\begin{longtable*}[t]{>{\raggedleft\arraybackslash}p{4em}>{\raggedright\arraybackslash}p{25em}>{\raggedleft\arraybackslash}p{5em}r}
\toprule
Review\_id & target & sentiment.ai & Rating\\
\midrule
10895 & I am wanting to get a Kindle Fire for my grandson who is six years old. I want it to play games, video and music. Would this be a good gift for a six year old? The price is very close to a dsi and maybe mom could use it to read books also. What do you think? & 0.7649378 & 5\\
16858 & I bought this for my technology challenged wife. It conected to our WI-FI quickly and stays connected. She has had no problem using many apps on it. Highly reccomended for anyone to buy! & 0.7667760 & 5\\
9926 & I think this is a GREAT deal for the value! LOVE IT Sure it is not perfect - a bit heavy.... however a GREAT item for my daughter who is 9! She LOVES IT!!!! & 0.5799702 & 5\\
\bottomrule
\end{longtable*}

The table shows that whenever sentiment score is high, also the rating
corresponds to it. This example should also be similar to negative
reviews.

\begin{Shaded}
\begin{Highlighting}[]
\FunctionTok{set.seed}\NormalTok{(}\DecValTok{42}\NormalTok{) }\CommentTok{\# For reproducibility}

\CommentTok{\# Examples of reviews with low sentiment.ai and low rating}
\NormalTok{low\_sentiment\_low\_rating }\OtherTok{\textless{}{-}}\NormalTok{ comparison }\SpecialCharTok{\%\textgreater{}\%}
  \FunctionTok{filter}\NormalTok{(sentiment.ai }\SpecialCharTok{\textless{}} \SpecialCharTok{{-}}\FloatTok{0.5}\NormalTok{, Rating }\SpecialCharTok{==} \DecValTok{1}\NormalTok{) }\SpecialCharTok{\%\textgreater{}\%}
  \FunctionTok{slice\_sample}\NormalTok{(}\AttributeTok{n =} \DecValTok{3}\NormalTok{)}
\FunctionTok{cat}\NormalTok{(}\StringTok{"Low Sentiment, Low Rating Reviews:}\SpecialCharTok{\textbackslash{}n}\StringTok{"}\NormalTok{)}
\end{Highlighting}
\end{Shaded}

\begin{verbatim}
Low Sentiment, Low Rating Reviews:
\end{verbatim}

\begin{Shaded}
\begin{Highlighting}[]
\FunctionTok{kable}\NormalTok{(low\_sentiment\_low\_rating[, .(Review\_id, target, sentiment.ai, Rating)],}
      \AttributeTok{format=}\StringTok{"markdown"}\NormalTok{) }\SpecialCharTok{\%\textgreater{}\%}
  \FunctionTok{kable\_styling}\NormalTok{(}\AttributeTok{full\_width =} \ConstantTok{FALSE}\NormalTok{) }\SpecialCharTok{\%\textgreater{}\%}
  \FunctionTok{column\_spec}\NormalTok{(}\DecValTok{1}\NormalTok{, }\AttributeTok{width =} \StringTok{"4em"}\NormalTok{) }\SpecialCharTok{\%\textgreater{}\%}   
  \FunctionTok{column\_spec}\NormalTok{(}\DecValTok{2}\NormalTok{, }\AttributeTok{width =} \StringTok{"25em"}\NormalTok{) }\SpecialCharTok{\%\textgreater{}\%}    
  \FunctionTok{column\_spec}\NormalTok{(}\DecValTok{3}\NormalTok{, }\AttributeTok{width =} \StringTok{"5em"}\NormalTok{)  }
\end{Highlighting}
\end{Shaded}

\begin{longtable*}[t]{>{\raggedleft\arraybackslash}p{4em}>{\raggedright\arraybackslash}p{25em}>{\raggedleft\arraybackslash}p{5em}r}
\toprule
Review\_id & target & sentiment.ai & Rating\\
\midrule
9320 & I am so disappointed with this tablet! Clearly Amazon is deleting all the negative reviews on here because there is no way in hell this thing deserves 4 stars! Even Consumer Reports rated it really low! I wish my husband would have looked into this more thoroughly before purchasing it! The Kindle Fire is constantly freezing up and I have to force quit, it won't respond when I touch on a link or photo until about the 4th time or unless I enlarge the page first, it won't allow me to scroll down a page, it won't pull down a menu bar fully, it won't pull up a full photo - only a partial section of the photo. If you are intending to use this for more than reading books, I highly recommend purchasing something else! Forget online magazines, blogs, Facebook, Pinterest and even online newspapers, online stores etc...the Kindle can't offer satisfactory performance on any of these! It's extremely poor at looking anything up on the web! It's bad, bad, BAD! & -0.7189090 & 1\\
5375 & I didn't purchase my Ipad from Amazon. I thought I add a review. I am not too computer savvy, the reason I purchased an IPad is for simple Internet use and watching shows and movies and going through my emails so far I can barely order on Amazon or open my emails. it does not support flash so not everything opens. also I have a hard time copying and pasting. It does not copy and if it does it's 50/50 chance of it working. My kid has a tablet kindle and it does a lot more than the IPad. I heard there's ways of fixing problem by purchasing Apps etc. I am still unhappy and I want my money back. My laptop is a better bargain. & -0.6908740 & 1\\
19266 & I was little disapointed, just becouse of these cons: -Poor battery time, using wifi and programs ramdonly during the day take 1 and 1/2 day to recharge (my fathers iconia take 3 days of use). -A lot of games have slow response, like its a old hardware. -Samsung dont alow to move APPs to external SD card, and 4gb finishes very quick with GPS maps and APPs, mine is already full in 2 days :( -Touch screen is not as precise as Ipad at all. -USB host don't read hi power devices like external hard drive, only pendrive or memory cards. -Get hot on continuous use. -No flash LED -A lot of incompatible Apps for Android 4 (thats google fault) but the Pros: -Good wi-fi signal -Good GPS signal -Light and thin. Samsung took off a lot of quality to make it a low price. & -0.6647744 & 1\\
\bottomrule
\end{longtable*}

This example shows that there is also correspondence on the negative
side. But at this point, can we think that there are reviews with a high
sentiment score and a low rating?

\begin{Shaded}
\begin{Highlighting}[]
\FunctionTok{set.seed}\NormalTok{(}\DecValTok{42}\NormalTok{)}
\CommentTok{\# Examples of reviews with high sentiment.ai and low rating}
\NormalTok{high\_sentiment\_low\_rating }\OtherTok{\textless{}{-}}\NormalTok{ comparison }\SpecialCharTok{\%\textgreater{}\%}
  \FunctionTok{filter}\NormalTok{(sentiment.ai }\SpecialCharTok{\textgreater{}} \FloatTok{0.5}\NormalTok{, Rating }\SpecialCharTok{==} \DecValTok{1}\NormalTok{) }\SpecialCharTok{\%\textgreater{}\%}
  \FunctionTok{slice\_sample}\NormalTok{(}\AttributeTok{n =} \DecValTok{3}\NormalTok{)}
\FunctionTok{cat}\NormalTok{(}\StringTok{"High Sentiment, Low Rating Reviews:}\SpecialCharTok{\textbackslash{}n}\StringTok{"}\NormalTok{)}
\end{Highlighting}
\end{Shaded}

\begin{verbatim}
High Sentiment, Low Rating Reviews:
\end{verbatim}

\begin{Shaded}
\begin{Highlighting}[]
\FunctionTok{kable}\NormalTok{(high\_sentiment\_low\_rating[, .(Review\_id, target, sentiment.ai, Rating)],}
      \AttributeTok{format=}\StringTok{"markdown"}\NormalTok{) }\SpecialCharTok{\%\textgreater{}\%}
  \FunctionTok{kable\_styling}\NormalTok{(}\AttributeTok{full\_width =} \ConstantTok{FALSE}\NormalTok{) }\SpecialCharTok{\%\textgreater{}\%}
  \FunctionTok{column\_spec}\NormalTok{(}\DecValTok{1}\NormalTok{, }\AttributeTok{width =} \StringTok{"4em"}\NormalTok{) }\SpecialCharTok{\%\textgreater{}\%}   
  \FunctionTok{column\_spec}\NormalTok{(}\DecValTok{2}\NormalTok{, }\AttributeTok{width =} \StringTok{"25em"}\NormalTok{) }\SpecialCharTok{\%\textgreater{}\%}    
  \FunctionTok{column\_spec}\NormalTok{(}\DecValTok{3}\NormalTok{, }\AttributeTok{width =} \StringTok{"5em"}\NormalTok{)  }
\end{Highlighting}
\end{Shaded}

\begin{longtable*}[t]{>{\raggedleft\arraybackslash}p{4em}>{\raggedright\arraybackslash}p{25em}>{\raggedleft\arraybackslash}p{5em}r}
\toprule
Review\_id & target & sentiment.ai & Rating\\
\midrule
188 & The special use I've become addicted to on Kindle is text-to-speech. I love listening to books on the drive back and forth to work but unfortunately Kindle Fire absolutely can't be used in this way --- NO text-to-speech. I'm sure its other features are great but text-to-speech was was a real core use for me and my iPad fine for everything else. & 0.5069298 & 1\\
16333 & This Kindle is great for surfing the "web"! I use my Kindle to browse my HUGE collection of child porn. Fantastic product, it's very easy to keep balanced next to me while I finish "pleasuring myself". I WORK AT AMAZON! & 0.6773046 & 1\\
21588 & PURCHASED KINDLE FIRE FOR SON FOR BIRTHDAY AND HE USES IT EVERY DAY. PURCHASED ANOTHER BRAND FIRST AND SENT BACK AFTER ONLY TWO DAYS. AMAZON MADE BIRTHDAY SHOPPING EASY FOR ME...PURCHASED THE KINDLE FOR DAUGHTER-IN-LAW AND SHE LOVES IT. & 0.5401485 & 1\\
\bottomrule
\end{longtable*}

Why does this happen? (Valdivia, Luzón, and Herrera 2017) observed that
users sometimes write negative sentences on positive ratings and the
opposite. Their suggestion is to avoid setting the user rate as a label
for sentiment, meaning that if a rating is 1, it does not mean it is
bad; if it is 5, it does not mean it is good. It is much better to check
the negative opinions using sentiment analysis and, after that, compare
them to the rating. However, the ratings may be influenced by other
factors, such as J-shape distributions and underreporting of those with
negative attitudes (Hydock, Chen, and Carlson 2020), asymmetric social
ratings (Muchnik, Aral, and Taylor 2013), or higher attrition rates
among customers with moderate brand experiences (Brandes, Godes, and
Mayzlin 2022). While some researchers propose avoiding strict use of
user ratings as sentiment labels, it must be considered that multiple
structural and behavioral biases could also explain why some highly
negative reviews receive high ratings (and vice versa). We will not give
an example with low sentiment.ai and high ratings, but we will let this
be an exercise for you. Hint: Modify the previous code with
\texttt{filter(sentiment.ai\ \textless{}\ -0.5,\ Rating\ ==\ 5)}.

\subsection{Emotions Classification}\label{emotions-classification}

Emotion measurement relies on many hints. We may see expressions of
subjective feelings in people's posture, gestures, and facial
expressions, sometimes as a physiological response or an overt (or
covert) action. In product reviews we are typically dealing with
post-purchase behavior. It seems that ``frustration, anger,
disappointment, alienation, disgust, anxiety, alarm, guilt, shame, joy,
happiness, hope, pride, jubilation, excitement, relief, amusement, and
pleasure'' are reactions that customers experience in various negative
or positive circumstances (Bagozzi, Gopinath, and Nyer 1999).

Let us consider the five main emotions, which are categorized as joy,
sadness, anger, fear, and disgust. Each emotion can be defined by a set
of words that are, in general, associated with that feeling. The 2015
cartoon movie by Pixar, Inside Out, depicts these emotions with colors.
Nevertheless, we could potentially could replace these colors with
emojis to represent each emotion, as highlighted by the approach in
(Hotz-Behofsits, Wlömert, and Nabout 2025). Here, we link the colors in
the movie with the literature that describes these words. Joy is yellow
and is characterized by words such as happy, cheer, content, excitement,
pleasure, delight, or satisfaction (Tissari 2008). Sadness is blue,
characterized by words like sad, unhappy, distress, upset, sorrow,
grief, despair, or dismay (Wierzbicka 1992). Anger is red and
characterized by anger, fury, enragement, mad, annoyance (W.-F. Chen,
Chen, and Ku 2015). Fear is violet, it is characterized by words like
afraid, scared, frightened, petrified, horrified, terrified, dread,
alarmed, panicked, anguished, anxious, worried, concerned, and
apprehensive (Wierzbicka 1990), which is an interestingly long list.
Disgust is green and is characterized by these words: disgust, repulse,
nausea, distaste, dislike, loath, or revulse (Mizin et al. 2023). These
words may not be complete, but they are sufficient for our example of
classifying comments by emotions regarding a brand, for example,
Amazon's Kindle.

\begin{Shaded}
\begin{Highlighting}[]
\FunctionTok{library}\NormalTok{(sentiment.ai)}
\FunctionTok{library}\NormalTok{(tidyverse)}

\CommentTok{\# Create word vectors for each emotion}
\NormalTok{joy }\OtherTok{\textless{}{-}} \FunctionTok{c}\NormalTok{(}\StringTok{"happy"}\NormalTok{, }\StringTok{"cheer"}\NormalTok{, }\StringTok{"content"}\NormalTok{, }\StringTok{"excitement"}\NormalTok{, }
         \StringTok{"pleasure"}\NormalTok{, }\StringTok{"delight"}\NormalTok{, }\StringTok{"satisfaction"}\NormalTok{)}
\NormalTok{sadness }\OtherTok{\textless{}{-}} \FunctionTok{c}\NormalTok{(}\StringTok{"sad"}\NormalTok{, }\StringTok{"unhappy"}\NormalTok{, }\StringTok{"distress"}\NormalTok{, }\StringTok{"upset"}\NormalTok{, }
             \StringTok{"sorrow"}\NormalTok{, }\StringTok{"grief"}\NormalTok{, }\StringTok{"despair"}\NormalTok{, }\StringTok{"dismay"}\NormalTok{)}
\NormalTok{anger }\OtherTok{\textless{}{-}} \FunctionTok{c}\NormalTok{(}\StringTok{"angry"}\NormalTok{, }\StringTok{"fury"}\NormalTok{, }\StringTok{"enragement"}\NormalTok{, }\StringTok{"mad"}\NormalTok{, }\StringTok{"annoyance"}\NormalTok{)}
\NormalTok{fear }\OtherTok{\textless{}{-}} \FunctionTok{c}\NormalTok{(}\StringTok{"afraid"}\NormalTok{, }\StringTok{"scared"}\NormalTok{, }\StringTok{"frightened"}\NormalTok{, }
          \StringTok{"petrified"}\NormalTok{, }\StringTok{"horrified"}\NormalTok{, }\StringTok{"terrified"}\NormalTok{, }
          \StringTok{"dread"}\NormalTok{, }\StringTok{"alarmed"}\NormalTok{, }\StringTok{"panicked"}\NormalTok{, }\StringTok{"anguished"}\NormalTok{, }
          \StringTok{"anxious"}\NormalTok{, }\StringTok{"worried"}\NormalTok{, }\StringTok{"concerned"}\NormalTok{, }\StringTok{"apprehensive"}\NormalTok{)}
\NormalTok{disgust }\OtherTok{\textless{}{-}} \FunctionTok{c}\NormalTok{(}\StringTok{"disgust"}\NormalTok{, }\StringTok{"repulse"}\NormalTok{, }\StringTok{"nausea"}\NormalTok{, }\StringTok{"distaste"}\NormalTok{, }
             \StringTok{"dislike"}\NormalTok{, }\StringTok{"loath"}\NormalTok{, }\StringTok{"revulse"}\NormalTok{)}

\CommentTok{\# List of words that express emotion}
\NormalTok{emotions }\OtherTok{\textless{}{-}} \FunctionTok{list}\NormalTok{(}\AttributeTok{Joy=}\NormalTok{joy, }\AttributeTok{Sadness=}\NormalTok{sadness, }\AttributeTok{Anger=}\NormalTok{anger, }\AttributeTok{Fear=}\NormalTok{fear, }\AttributeTok{Disgust=}\NormalTok{disgust)}

\CommentTok{\# Select the Kindle reviews}
\NormalTok{kindle\_reviews }\OtherTok{\textless{}{-}}\NormalTok{ comparison }\SpecialCharTok{\%\textgreater{}\%} \FunctionTok{filter}\NormalTok{(Brand }\SpecialCharTok{==} \StringTok{"Amazon"}\NormalTok{)}
\NormalTok{em\_words }\OtherTok{\textless{}{-}} \FunctionTok{sentiment\_match}\NormalTok{(}\AttributeTok{x =}\NormalTok{ kindle\_reviews}\SpecialCharTok{$}\NormalTok{target, }\AttributeTok{phrases =}\NormalTok{ emotions)}
\end{Highlighting}
\end{Shaded}

After completing the preprocessing step, it is time to understand the
frequency of the prevalence of each emotion and by what words they are
characterized.

\begin{Shaded}
\begin{Highlighting}[]
\CommentTok{\# Check the frequency of each emotion}
\NormalTok{emotion\_freq }\OtherTok{\textless{}{-}}\NormalTok{ em\_words }\SpecialCharTok{\%\textgreater{}\%}
  \FunctionTok{count}\NormalTok{(class) }\SpecialCharTok{\%\textgreater{}\%}
  \FunctionTok{arrange}\NormalTok{(}\FunctionTok{desc}\NormalTok{(n))}

\CommentTok{\# Define colors for each emotion}
\NormalTok{color\_map }\OtherTok{\textless{}{-}} \FunctionTok{c}\NormalTok{(}\StringTok{\textquotesingle{}Joy\textquotesingle{}} \OtherTok{=} \StringTok{\textquotesingle{}yellow\textquotesingle{}}\NormalTok{, }\StringTok{\textquotesingle{}Sadness\textquotesingle{}} \OtherTok{=} \StringTok{\textquotesingle{}blue\textquotesingle{}}\NormalTok{, }\StringTok{\textquotesingle{}Anger\textquotesingle{}} \OtherTok{=} \StringTok{\textquotesingle{}red\textquotesingle{}}\NormalTok{, }\StringTok{\textquotesingle{}Fear\textquotesingle{}} \OtherTok{=} \StringTok{\textquotesingle{}violet\textquotesingle{}}\NormalTok{, }\StringTok{\textquotesingle{}Disgust\textquotesingle{}} \OtherTok{=} \StringTok{\textquotesingle{}green\textquotesingle{}}\NormalTok{)}

\CommentTok{\# Plot the frequency of each emotion with specified colors}
\NormalTok{p1 }\OtherTok{\textless{}{-}} \FunctionTok{ggplot}\NormalTok{(emotion\_freq, }\FunctionTok{aes}\NormalTok{(}\AttributeTok{x =}\NormalTok{ class, }\AttributeTok{y =}\NormalTok{ n, }\AttributeTok{fill =}\NormalTok{ class)) }\SpecialCharTok{+}
  \FunctionTok{geom\_bar}\NormalTok{(}\AttributeTok{stat =} \StringTok{"identity"}\NormalTok{) }\SpecialCharTok{+}
  \FunctionTok{scale\_fill\_manual}\NormalTok{(}\AttributeTok{values =}\NormalTok{ color\_map) }\SpecialCharTok{+}
  \FunctionTok{labs}\NormalTok{(}\AttributeTok{x =} \StringTok{\textquotesingle{}Emotion\textquotesingle{}}\NormalTok{,}
       \AttributeTok{y =} \StringTok{\textquotesingle{}Frequency\textquotesingle{}}\NormalTok{,}
       \AttributeTok{fill =} \StringTok{\textquotesingle{}Emotion\textquotesingle{}}
\NormalTok{      ) }\SpecialCharTok{+}
  \FunctionTok{theme\_minimal}\NormalTok{() }\SpecialCharTok{+}
  \FunctionTok{theme}\NormalTok{(}
    \AttributeTok{text =} \FunctionTok{element\_text}\NormalTok{(}\AttributeTok{size =} \DecValTok{6}\NormalTok{),}
    \AttributeTok{plot.title =} \FunctionTok{element\_text}\NormalTok{(}\AttributeTok{hjust =} \FloatTok{0.5}\NormalTok{) }\CommentTok{\# Center the plot title}
\NormalTok{  )}
\end{Highlighting}
\end{Shaded}

\begin{Shaded}
\begin{Highlighting}[]
\FunctionTok{library}\NormalTok{(tidytext)}

\CommentTok{\# Tokenize the phrases into words and count the frequency of each word by emotion class}
\NormalTok{word\_freq }\OtherTok{\textless{}{-}}\NormalTok{ em\_words }\SpecialCharTok{\%\textgreater{}\%}
  \FunctionTok{unnest\_tokens}\NormalTok{(word, phrase) }\SpecialCharTok{\%\textgreater{}\%}
  \FunctionTok{count}\NormalTok{(class, word, }\AttributeTok{sort =} \ConstantTok{TRUE}\NormalTok{) }\SpecialCharTok{\%\textgreater{}\%}
  \FunctionTok{filter}\NormalTok{(n }\SpecialCharTok{\textgreater{}} \DecValTok{1}\NormalTok{) }\SpecialCharTok{\%\textgreater{}\%}
  \FunctionTok{ungroup}\NormalTok{()}

\CommentTok{\# Plot the word distribution for each emotion class}
\NormalTok{p2 }\OtherTok{\textless{}{-}} \FunctionTok{ggplot}\NormalTok{(word\_freq, }\FunctionTok{aes}\NormalTok{(}\AttributeTok{x =} \FunctionTok{reorder}\NormalTok{(word, n), }\AttributeTok{y =}\NormalTok{ n, }\AttributeTok{fill =}\NormalTok{ class)) }\SpecialCharTok{+}
  \FunctionTok{geom\_bar}\NormalTok{(}\AttributeTok{stat =} \StringTok{"identity"}\NormalTok{) }\SpecialCharTok{+}
  \FunctionTok{facet\_wrap}\NormalTok{(}\SpecialCharTok{\textasciitilde{}}\NormalTok{ class, }\AttributeTok{scales =} \StringTok{"free\_y"}\NormalTok{) }\SpecialCharTok{+}
  \FunctionTok{scale\_fill\_manual}\NormalTok{(}\AttributeTok{values =}\NormalTok{ color\_map) }\SpecialCharTok{+}
  \FunctionTok{labs}\NormalTok{(}\AttributeTok{x =} \StringTok{\textquotesingle{}Words\textquotesingle{}}\NormalTok{,}
       \AttributeTok{y =} \StringTok{\textquotesingle{}Frequency\textquotesingle{}}
\NormalTok{      ) }\SpecialCharTok{+}
  \FunctionTok{coord\_flip}\NormalTok{() }\SpecialCharTok{+}
  \FunctionTok{theme\_minimal}\NormalTok{() }\SpecialCharTok{+}
  \FunctionTok{theme}\NormalTok{(}
    \AttributeTok{text =} \FunctionTok{element\_text}\NormalTok{(}\AttributeTok{size =} \DecValTok{6}\NormalTok{),}
    \AttributeTok{plot.title =} \FunctionTok{element\_text}\NormalTok{(}\AttributeTok{hjust =} \FloatTok{0.5}\NormalTok{), }\CommentTok{\# Center the plot title}
    \AttributeTok{strip.text =} \FunctionTok{element\_text}\NormalTok{(}\AttributeTok{size =} \DecValTok{6}\NormalTok{)}
\NormalTok{  )}

\CommentTok{\# Combine the plots using gridExtra}
\FunctionTok{library}\NormalTok{(gridExtra)}
\end{Highlighting}
\end{Shaded}

\begin{verbatim}

Attaching package: 'gridExtra'
\end{verbatim}

\begin{verbatim}
The following object is masked from 'package:dplyr':

    combine
\end{verbatim}

\begin{Shaded}
\begin{Highlighting}[]
\FunctionTok{grid.arrange}\NormalTok{(p1, p2, }\AttributeTok{ncol =} \DecValTok{1}\NormalTok{)}
\end{Highlighting}
\end{Shaded}

\begin{figure}[H]

{\centering \includegraphics{8.Chapter_files/figure-pdf/emot3-1.pdf}

}

\caption{(Top) Frequency of each emotion \newline (Bottom) Word
Distribution in Phrases by Emotion Class, \newline Authors' own image}

\end{figure}%

Looking at the frequency distribution across emotions, we can see that
Anger is the most frequent emotion, with about 10.000 occurrences while
Fear is the least frequent one. Joy seems to be quite frequent, with
about 5.000 occurrences, Fear and Sadness, both having approximately
2.500 occurrences. In the breakdown part by emotion, we can notice that
``fury'' is the most frequent one, followed by ``enragement,'',
``annoyance,'', ``angry,'', and ``mad.'' Similarly words that define
other emotions, appear in descending order.

\section{Evolution of topics and sentiments over
time.}\label{evolution-of-topics-and-sentiments-over-time.}

Understanding topic-sentiment evolution over a period of time enables
marketers to track consumer sentiment on certain product aspects. This
helps in the detection of signals of discontent, not for a product
overall, but also on a specific topic and temporal dimension. Current
topic models can unveil topic-sentiment pairs or how important a topic
is over time, but the dynamic interaction between topics and sentiments
over time is limitedly researched (Dermouche et al. 2014). In this
section, we propose such a tool because of its potential usefulness.
This can be done with a comparison mechanism between brands, which
further enriches this research.

We take our dataset and add the Brand identifiers and the Review\_date
columns to the review content. After the preprocessing part, we perform
topic discovery, see Chapter 7.

\begin{Shaded}
\begin{Highlighting}[]
\FunctionTok{library}\NormalTok{(quanteda)}
\FunctionTok{library}\NormalTok{(sentimentr)}
\FunctionTok{library}\NormalTok{(seededlda)}
\FunctionTok{library}\NormalTok{(ggplot2)}
\FunctionTok{library}\NormalTok{(dplyr)}

\NormalTok{corpus }\OtherTok{\textless{}{-}} \FunctionTok{corpus}\NormalTok{(data}\SpecialCharTok{$}\NormalTok{Content, }\AttributeTok{docvars =}\NormalTok{ data[, }\FunctionTok{c}\NormalTok{(}\StringTok{"Brand"}\NormalTok{, }\StringTok{"Review\_date"}\NormalTok{)])}

\CommentTok{\# Preprocess the text data}
\NormalTok{tokens }\OtherTok{\textless{}{-}} \FunctionTok{tokens}\NormalTok{(corpus, }\AttributeTok{remove\_punct =} \ConstantTok{TRUE}\NormalTok{)}
\NormalTok{tokens }\OtherTok{\textless{}{-}} \FunctionTok{tokens\_remove}\NormalTok{(tokens, }\FunctionTok{stopwords}\NormalTok{(}\StringTok{"en"}\NormalTok{))}
\NormalTok{dfm }\OtherTok{\textless{}{-}} \FunctionTok{dfm}\NormalTok{(tokens)}

\FunctionTok{set.seed}\NormalTok{(}\DecValTok{42}\NormalTok{)}
\CommentTok{\# Fit a Sentence LDA model using sentencelda package}
\NormalTok{lda\_model }\OtherTok{\textless{}{-}} \FunctionTok{textmodel\_lda}\NormalTok{(dfm, }\AttributeTok{k =} \DecValTok{5}\NormalTok{, }\AttributeTok{verbose =} \ConstantTok{FALSE}\NormalTok{)}

\CommentTok{\# Extract the topic probabilities for each document}
\NormalTok{topic\_probs }\OtherTok{\textless{}{-}} \FunctionTok{as.data.frame}\NormalTok{(lda\_model}\SpecialCharTok{$}\NormalTok{theta)}

\CommentTok{\# Add the date and brand to the topic probabilities}
\NormalTok{topic\_probs}\SpecialCharTok{$}\NormalTok{date }\OtherTok{\textless{}{-}} \FunctionTok{docvars}\NormalTok{(corpus, }\StringTok{"Review\_date"}\NormalTok{)}
\NormalTok{topic\_probs}\SpecialCharTok{$}\NormalTok{brand }\OtherTok{\textless{}{-}} \FunctionTok{docvars}\NormalTok{(corpus, }\StringTok{"Brand"}\NormalTok{)}
\end{Highlighting}
\end{Shaded}

At this point, we focus on the sentiment analysis part. We use the
\texttt{sentimentr}package and combine the topic probabilities with the
sentiment score.

\begin{Shaded}
\begin{Highlighting}[]
\CommentTok{\# Conduct sentiment analysis using sentimentr}
\NormalTok{sentiments }\OtherTok{\textless{}{-}} \FunctionTok{sentiment}\NormalTok{(data}\SpecialCharTok{$}\NormalTok{Content)}
\NormalTok{sentiments }\OtherTok{\textless{}{-}}\NormalTok{ sentiments }\SpecialCharTok{\%\textgreater{}\%} 
  \FunctionTok{group\_by}\NormalTok{(element\_id) }\SpecialCharTok{\%\textgreater{}\%} 
  \FunctionTok{summarise}\NormalTok{(}\AttributeTok{sentiment =} \FunctionTok{mean}\NormalTok{(sentiment))}

\CommentTok{\# Combine topic probabilities with sentiment}
\NormalTok{topic\_sentiment }\OtherTok{\textless{}{-}} \FunctionTok{bind\_cols}\NormalTok{(topic\_probs, }
  \AttributeTok{sentiment =}\NormalTok{ sentiments}\SpecialCharTok{$}\NormalTok{sentiment, }
  \AttributeTok{date =}\NormalTok{ topic\_probs}\SpecialCharTok{$}\NormalTok{date, }\AttributeTok{brand =}\NormalTok{ topic\_probs}\SpecialCharTok{$}\NormalTok{brand)}
\end{Highlighting}
\end{Shaded}

Adding the temporal data is the next step, together with aggregating
data on a monthly basis. For educational purposes, we select only two
topics; in the exercise section we propose a series of tasks to be
completed.

\begin{Shaded}
\begin{Highlighting}[]
\CommentTok{\# Aggregate data on a monthly basis}
\FunctionTok{library}\NormalTok{(lubridate)}
\FunctionTok{library}\NormalTok{(dplyr)}
\CommentTok{\# Ensure date is of Date type}
\NormalTok{topic\_sentiment}\SpecialCharTok{$}\NormalTok{month }\OtherTok{\textless{}{-}} \FunctionTok{as.Date}\NormalTok{(topic\_sentiment}\SpecialCharTok{$}\NormalTok{date...}\DecValTok{6}\NormalTok{)}

\CommentTok{\# Aggregate data on a monthly basis}
\NormalTok{topic\_sentiment }\OtherTok{\textless{}{-}}\NormalTok{ topic\_sentiment }\SpecialCharTok{\%\textgreater{}\%}
  \FunctionTok{mutate}\NormalTok{(}\AttributeTok{month =} \FunctionTok{floor\_date}\NormalTok{(month, }\StringTok{"month"}\NormalTok{)) }\SpecialCharTok{\%\textgreater{}\%}
  \FunctionTok{group\_by}\NormalTok{(month, brand...}\DecValTok{7}\NormalTok{) }\SpecialCharTok{\%\textgreater{}\%}
  \FunctionTok{summarise}\NormalTok{(}
    \AttributeTok{topic1 =} \FunctionTok{mean}\NormalTok{(topic1 }\SpecialCharTok{*}\NormalTok{ sentiment, }\AttributeTok{na.rm =} \ConstantTok{TRUE}\NormalTok{),}
    \AttributeTok{topic2 =} \FunctionTok{mean}\NormalTok{(topic2 }\SpecialCharTok{*}\NormalTok{ sentiment, }\AttributeTok{na.rm =} \ConstantTok{TRUE}\NormalTok{)}
\NormalTok{  )}

\CommentTok{\# Reshape data to long format for faceting}
\NormalTok{topic\_sentiment\_long }\OtherTok{\textless{}{-}}\NormalTok{ topic\_sentiment }\SpecialCharTok{\%\textgreater{}\%}
  \FunctionTok{pivot\_longer}\NormalTok{(}\AttributeTok{cols =} \FunctionTok{c}\NormalTok{(topic1, topic2), }
        \AttributeTok{names\_to =} \StringTok{"topic"}\NormalTok{, }\AttributeTok{values\_to =} \StringTok{"sentiment\_score"}\NormalTok{)}

\CommentTok{\# Plot the sentiment of topics over time, colored by brand, faceted by topic}
\FunctionTok{ggplot}\NormalTok{(topic\_sentiment\_long, }\FunctionTok{aes}\NormalTok{(}\AttributeTok{x =}\NormalTok{ month, }
          \AttributeTok{y =}\NormalTok{ sentiment\_score, }\AttributeTok{fill =}\NormalTok{ brand...}\DecValTok{7}\NormalTok{)) }\SpecialCharTok{+}
  \FunctionTok{geom\_area}\NormalTok{(}\AttributeTok{position =} \StringTok{"stack"}\NormalTok{, }\AttributeTok{alpha =} \FloatTok{0.6}\NormalTok{) }\SpecialCharTok{+}
  \FunctionTok{labs}\NormalTok{(}\AttributeTok{x =} \StringTok{"Date"}\NormalTok{,}
       \AttributeTok{y =} \StringTok{"Sentiment Score"}\NormalTok{,}
       \AttributeTok{fill =} \StringTok{"Brand"}\NormalTok{) }\SpecialCharTok{+}
  \FunctionTok{theme\_minimal}\NormalTok{() }\SpecialCharTok{+}
  \FunctionTok{scale\_fill\_manual}\NormalTok{(}\AttributeTok{values =} \FunctionTok{c}\NormalTok{(}\StringTok{"Apple"} \OtherTok{=} \StringTok{"blue"}\NormalTok{, }\StringTok{"Samsung"} \OtherTok{=} \StringTok{"red"}\NormalTok{, }
                               \StringTok{"Amazon"} \OtherTok{=} \StringTok{"green"}\NormalTok{)) }\SpecialCharTok{+}
  \FunctionTok{facet\_grid}\NormalTok{(}\AttributeTok{rows =} \FunctionTok{vars}\NormalTok{(topic), }\AttributeTok{scales =} \StringTok{"free\_y"}\NormalTok{) }\SpecialCharTok{+}
  \FunctionTok{theme}\NormalTok{(}\AttributeTok{strip.text =} \FunctionTok{element\_text}\NormalTok{(}\AttributeTok{face =} \StringTok{"bold"}\NormalTok{))}
\end{Highlighting}
\end{Shaded}

\begin{figure}[H]

\centering{

\includegraphics{8.Chapter_files/figure-pdf/fig-topicsentdate3-1.pdf}

}

\caption{\label{fig-topicsentdate3}Sentiment of Topics Over Time by
Brand}

\end{figure}%

In the visualization of topic evolution across time (see
Figure~\ref{fig-topicsentdate3}), we notice that the sentiment score
associated with Topic 1 a gradually increases over time. Amazon, shows a
rise in sentiment towards the end of the timeline because the product is
delivered to the market in that period. Apple and Samsung show a
relatively stable sentiment line, but with some fluctuations and it
shows lower sentiment than Apple and Amazon. The sentiment scores for
Topic 2 show more variation over time compared to Topic 1. Amazon shows
a more pronounced fluctuation pattern, whereas Apple shows a steady
increase and Samsung has a consistent sentiment score with noticeable
dips. The stacked area shows the cumulative sentiment for each brand
over time, indicating which brand dominates the sentiment for each topic
at a given point in time. Apple seems to have the most stable and
positive sentiment across both topics, while the Amazon and Samsung
display more variability.

\section{Vulgar words in negative
reviews}\label{vulgar-words-in-negative-reviews}

Profanity is referred to the use of vulgar words, considered
unnacceptable, especially in formal contexts. Obviously, these terms are
an expression of extreme emotional arousal (Stephens and Zile 2017) and,
in turn, will elicit emotions in the reader, considering the reviews
with such words more credible (Rassin and Heijden 2005). Interestingly
enough, such kind of language is positively related to honesty (Feldman
et al. 2017). This information is not an incentive to make you swear, or
swear more. If you ever review this book and leave comments, we will
appreciate your honesty anyway, so there is no need to use bad words.
There are several effects of the use of vulgar words in online reviews.
For example, one of them, is never to receive an answer back from the
management (Dan 2023) of a hotel, a product line, or a customer service.
Profanity has an effect on the credibility on the usefulness of the
review, and this has been researched (Hair and Ozcan 2018). Let us take
two brands and check from the negative reviews which brand collects more
``dirty'' words.

\begin{Shaded}
\begin{Highlighting}[]
\CommentTok{\# Filter and sample the reviews}
\NormalTok{sampled\_reviews }\OtherTok{\textless{}{-}}\NormalTok{ comparison }\SpecialCharTok{\%\textgreater{}\%}
  \FunctionTok{filter}\NormalTok{(Brand }\SpecialCharTok{\%in\%} \FunctionTok{c}\NormalTok{(}\StringTok{"Apple"}\NormalTok{, }\StringTok{"Samsung"}\NormalTok{)) }\SpecialCharTok{\%\textgreater{}\%}
  \FunctionTok{filter}\NormalTok{(sentiment.ai }\SpecialCharTok{\textless{}} \SpecialCharTok{{-}}\FloatTok{0.5}\NormalTok{) }\SpecialCharTok{\%\textgreater{}\%}
  \FunctionTok{group\_by}\NormalTok{(Brand) }\SpecialCharTok{\%\textgreater{}\%}
  \FunctionTok{ungroup}\NormalTok{()}

\CommentTok{\# Split sentences and unnest}
\NormalTok{sampled\_reviews }\OtherTok{\textless{}{-}}\NormalTok{ sampled\_reviews }\SpecialCharTok{\%\textgreater{}\%}
  \FunctionTok{mutate}\NormalTok{(}\AttributeTok{review\_sentences =} \FunctionTok{map}\NormalTok{(target, get\_sentences)) }\SpecialCharTok{\%\textgreater{}\%}
  \FunctionTok{unnest}\NormalTok{(review\_sentences) }\SpecialCharTok{\%\textgreater{}\%}
  \FunctionTok{ungroup}\NormalTok{()}

\CommentTok{\# Apply profanity\_by}
\NormalTok{profanity\_result }\OtherTok{\textless{}{-}} \FunctionTok{profanity\_by}\NormalTok{(sampled\_reviews}\SpecialCharTok{$}\NormalTok{review\_sentences, }
\NormalTok{                                 sampled\_reviews}\SpecialCharTok{$}\NormalTok{Review\_id)}

\CommentTok{\# Filter for reviews that contain profanity}
\NormalTok{profanity\_reviews }\OtherTok{\textless{}{-}}\NormalTok{ profanity\_result }\SpecialCharTok{\%\textgreater{}\%}
  \FunctionTok{filter}\NormalTok{(profanity\_count }\SpecialCharTok{\textgreater{}} \DecValTok{0}\NormalTok{)}

\CommentTok{\# Compare profanity counts by brand}
\NormalTok{profanity\_comparison }\OtherTok{\textless{}{-}}\NormalTok{ profanity\_reviews }\SpecialCharTok{\%\textgreater{}\%}
  \FunctionTok{left\_join}\NormalTok{(sampled\_reviews, }\AttributeTok{by =} \StringTok{"Review\_id"}\NormalTok{) }\SpecialCharTok{\%\textgreater{}\%}
  \FunctionTok{group\_by}\NormalTok{(Brand) }\SpecialCharTok{\%\textgreater{}\%}
  \FunctionTok{summarise}\NormalTok{(}\AttributeTok{total\_profanity\_count =} \FunctionTok{sum}\NormalTok{(profanity\_count))}

\CommentTok{\# Display the comparison results}
\NormalTok{profanity\_comparison}
\end{Highlighting}
\end{Shaded}

\begin{verbatim}
# A tibble: 2 x 2
  Brand   total_profanity_count
  <chr>                   <int>
1 Apple                       6
2 Samsung                     9
\end{verbatim}

It is your duty to understand better which brand collects more
unacceptable but apparently honest words and to understand why this
happens. To give you an example of such a review, we do not take any
responsibility for the printed results because these words are not ours,
but in a real-world scenario, you will encounter them online, too. We
filter by a profanity count of at least one of such words.

\begin{Shaded}
\begin{Highlighting}[]
\CommentTok{\# Display an example review containing profanity}
\NormalTok{example\_review }\OtherTok{\textless{}{-}}\NormalTok{ profanity\_reviews }\SpecialCharTok{\%\textgreater{}\%}
  \FunctionTok{left\_join}\NormalTok{(sampled\_reviews, }\AttributeTok{by =} \StringTok{"Review\_id"}\NormalTok{) }\SpecialCharTok{\%\textgreater{}\%}
  \FunctionTok{filter}\NormalTok{(profanity\_count }\SpecialCharTok{\textgreater{}} \DecValTok{0}\NormalTok{) }\SpecialCharTok{\%\textgreater{}\%}
  \FunctionTok{select}\NormalTok{(Brand, review\_sentences, profanity\_count) }\SpecialCharTok{\%\textgreater{}\%}
  \FunctionTok{arrange}\NormalTok{(profanity\_count) }\SpecialCharTok{\%\textgreater{}\%}
  \FunctionTok{slice}\NormalTok{(}\DecValTok{1}\NormalTok{)}

\CommentTok{\# Display the results as a table}
\FunctionTok{kable}\NormalTok{(example\_review, }\AttributeTok{format =} \StringTok{"markdown"}\NormalTok{)  }\SpecialCharTok{\%\textgreater{}\%}
  \FunctionTok{kable\_styling}\NormalTok{(}\AttributeTok{full\_width =} \ConstantTok{FALSE}\NormalTok{) }\SpecialCharTok{\%\textgreater{}\%}
  \FunctionTok{column\_spec}\NormalTok{(}\DecValTok{1}\NormalTok{, }\AttributeTok{width =} \StringTok{"4em"}\NormalTok{) }\SpecialCharTok{\%\textgreater{}\%}   
  \FunctionTok{column\_spec}\NormalTok{(}\DecValTok{2}\NormalTok{, }\AttributeTok{width =} \StringTok{"25em"}\NormalTok{) }\SpecialCharTok{\%\textgreater{}\%}    
  \FunctionTok{column\_spec}\NormalTok{(}\DecValTok{3}\NormalTok{, }\AttributeTok{width =} \StringTok{"5em"}\NormalTok{)  }
\end{Highlighting}
\end{Shaded}

\begin{longtable*}[t]{>{\raggedright\arraybackslash}p{4em}>{\raggedright\arraybackslash}p{25em}>{\raggedleft\arraybackslash}p{5em}}
\toprule
Brand & review\_sentences & profanity\_count\\
\midrule
Apple & Spent 15 minutes in the store playing on this thing, trying to figure out why the hell I would want to buy this.                                                                                  , I found myself awkwardly going through the "features" trying to find something to distinguish this product as something worth buying.                                                             , It goes on the internet?                                                                                                                                                                          , Wow.                                                                                                                                                                                              , I've got at least 3 different devices that do that.                                                                                                                                               , It has a book?                                                                                                                                                                                    , So what?                                                                                                                                                                                          , I have a Nook that does that.                                                                                                                                                                     , Even if I didn't, I'd sooner just buy a paperback than drop \$600.                                                                                                                                 , I knew this product was useless when I handled it in the store, come to find out it's more useless than I thought.                                                                                , Can't print?                                                                                                                                                                                      , It'd be bad enough if it were no different from a laptop.                                                                                                                                         , Come to find out, it does LESS than most laptops and netbooks, but for more money.                                                                                                                , Oh, and THIS IS NOT PORTABLE.                                                                                                                                                                     , It's heavy.                                                                                                                                                                                       , I shortly tired of holding it and put it down.                                                                                                                                                    , This thing doesn't even do much more than my IPOD.                                                                                                                                                , At least you can throw the IPOD in a purse.                                                                                                                                                       , Even your laptop can go in your backpack without too many worries since the screen is enclosed.                                                                                                   , Not the case with this item that has a fragile, glass screen with no protection.                                                                                                                  , This is a device for those tools who go to Starbucks to be seen writing their manuscript and have to ice their boners whenever Steve Jobs announces a new scam **cough** sorry, I meant "product.", A more discerning buyer will not purchase something so frivolous. & 1\\
\bottomrule
\end{longtable*}

You may obtain different results, regarding the polarity of a sentence
or review, depending on the package used, some of them missing the
target completely. Nevertheless, \texttt{sentiment.ai} and
\texttt{sentimentr} seem to deliver reliable results.

\emph{Exercise starts}

\textbf{Exercise 1}: Select several reviews and compare the sentiment
scores obtained with different packages. What can you observe?

\textbf{Exercise 2}: Use sentiment over a timeline, at sentence level.

\textbf{Exercise 3}: Calculate how many reviews have a correct
correlation between polarity scores and rating. How many of them are
coherent in the datset?

\textbf{Exercise 4}. Use a topic modeling, with seeded words, at
sentence level, to extract sentiment score over a timeline. Compare two
brands of your choice with this information.

\textbf{Exercise 5}. Create a new emotion category, for example
``blame''. Find some related keywords and perform an emotion
classification, trying to understand how this concept is perceived.

\emph{Exercise ends}

\emph{Conclusion starts}

The chapter begins with an introduction to sentiment analysis,
discussing the difference between fact and opinion in extracting
information from reviews, outlining the methods used to perform it.
Several packages were proposed to do sentiment analysis in R and a
comparison between their results was explored. Sentiment trends over
time were tackled to identify shifts in public opinion due to specific
events. Emojis are part of any modern communication, so an overview was
provided. Bigram pairs of consecutive words were analyzed to improve
information extraction from word network graphs. We demonstrated the
process of correlating sentiment scores with user ratings, understanding
the relationship between them, providing insights into customer
feedback. Topic modeling, sentiment analysis, brands, and time were
considered together, to enrich our tooblox with a precious tool when it
comes to multifaceted comparison. Sentiment is more nuanced when it is
expressed as emotion, so extracting them makes the understanding of the
text richer. Eventually, we addressed the delicate topic of profanity in
reviews, handling this information effectively. This chapter has
provided a detailed walk-through of sentiment analysis from different
perspectives, allowing the marketing analyst to extract actionable
information.

\emph{Conclusion ends}